\documentclass[12pt,a4paper,oneside]{report}

% --- Packages ---
\usepackage[utf8]{inputenc}
\usepackage[T1]{fontenc}
\usepackage{lmodern}
\usepackage{amsmath,amssymb,amsfonts}
\usepackage{graphicx}
\usepackage{geometry}
\usepackage{setspace}
\usepackage{booktabs}
\usepackage{array}
\usepackage{float}
\usepackage{caption}
\usepackage{subcaption}
\usepackage{xcolor}
\usepackage{enumitem}
\usepackage{pdfpages}
\usepackage{fancyhdr}
\usepackage{microtype}
\usepackage{svg}
\usepackage[colorlinks=true, linkcolor=blue, citecolor=blue, urlcolor=blue]{hyperref}

% --- Unicode Character Binds ---
\DeclareUnicodeCharacter{2013}{--}
\DeclareUnicodeCharacter{2014}{---}
\DeclareUnicodeCharacter{00D7}{\ensuremath{\times}}
\DeclareUnicodeCharacter{223C}{\ensuremath{\sim}}
\DeclareUnicodeCharacter{2248}{\ensuremath{\approx}}
\DeclareUnicodeCharacter{2192}{\ensuremath{\rightarrow}}
\DeclareUnicodeCharacter{00B0}{\ensuremath{^\circ}}

% --- Page Setup (GOST 2.105-1995 compliant) ---
\geometry{left=30mm, right=15mm, top=20mm, bottom=20mm}
\onehalfspacing
\setlength{\parindent}{1.25cm}
\setlength{\parskip}{0pt}
\graphicspath{{figures/}}
\raggedbottom

% --- Page Numbering (GOST 7.32-2001) ---
\pagestyle{fancy}
\fancyhf{}
\fancyfoot[C]{\thepage}
\renewcommand{\headrulewidth}{0pt}
\renewcommand{\footrulewidth}{0pt}

% --- Title Page ---
\title{
    \vspace{2cm}
    \textbf{\Large Adaptive Design of Embodied Robotic Systems for Real-Time Self-Modeling Toward Swarm Coordination} \\
    \vspace{1cm}
    \large Master's Thesis
}
\author{
    Muhammad Zeeshan Asghar \\
    Master's Program in Data Science \\
    HSE University
}
\date{\vspace{1cm} 2026}

\begin{document}

% --- Front Matter ---
\pagenumbering{gobble}
\maketitle
\clearpage

% --- Abstract ---
\chapter*{Abstract}
\addcontentsline{toc}{chapter}{Abstract}

Embodied robotic systems deployed in unstructured environments inevitably experience embodiment drift---the progressive divergence between a robot's internal kinematic model and its physical hardware caused by mechanical wear, calibration slip, thermal deformation, and unmodeled impacts \cite{bongard2006resilient, cully2015robots}. When a robot's internal model no longer matches its physical embodiment, traditional Jacobian-based inverse kinematics controllers fail catastrophically: end-effectors diverge from targets, oscillations destabilize the system, and contact-rich tasks generate unsafe forces \cite{kwiatkowski2019task, chen2022full}. Maintaining true autonomy requires self-modeling---the ability to infer one's own morphology directly from sensorimotor experience.

Contemporary visual self-modeling architectures based on Neural Radiance Fields (NeRF) and 3D Gaussian Splatting (3DGS) achieve high morphological expressiveness but suffer from a critical latency paradox \cite{hu2024teaching, hu2025learning}. A diagnostic self-model is only valuable if it can provide state estimates fast enough to influence the feedback controller. Modern robotic control loops commonly operate at 1 kHz, permitting a maximum latency of 1.0 millisecond per inference \cite{siciliano2009robotics}. Existing neural rendering self-models operate at 5--37 FPS---latencies of 27--192 milliseconds---rendering them unusable for real-time control integration. This represents a fundamental architectural mismatch: volumetric reconstruction is computationally incompatible with sub-millisecond control deadlines.

This thesis resolves this bottleneck by introducing NeuroKin, a lightweight fully convolutional direct sensorimotor decoder that completely bypasses explicit 3D volumetric reconstruction. NeuroKin maps joint angles directly to visual silhouettes, achieving 21.88 dB PSNR at 7,400 FPS (0.135 ms latency)---a 1,418$\times$ speedup over implicit NeRF baselines with superior fidelity \cite{hu2024teaching}. Training completes in 33.5 seconds on a Tesla T4 GPU. An auxiliary multi-task variant, ResNeuroKin-D, recovers geometric depth while maintaining 2,500 FPS (0.400 ms latency). Integration of NeuroKin's spatial estimates into a Jacobian-based feedback loop enables dynamic error recovery after severe fault injection, demonstrating that sub-millisecond visual self-modeling can actively stabilize a damaged manipulator.

Building on this individual diagnostic capability, we scale to collective fault tolerance via temporal stigmergy---a mechanism of indirect coordination wherein agents leave environmental traces that stimulate adaptive responses in subsequent agents. In a 100-stage sequential swarm under progressive fault injection, the standard inverse-kinematics baseline collapses from 100\% pre-fault success to 13\% post-fault success. Conversely, NeuroKin-enabled agents read and write a shared failure ledger, autonomously trigger localized recovery modes, and preserve 88\% systemic success.

This thesis provides four primary contributions: (1) formal benchmarking of the latency paradox in neural rendering self-models, (2) the NeuroKin direct decoding architecture achieving sub-millisecond inference, (3) visual closed-loop damage adaptation, and (4) empirical validation that fast individual self-modeling enables decentralized stigmergic coordination---the computational prerequisite for physical swarm deployment.

\clearpage

% --- Acknowledgements ---
\chapter*{Acknowledgements}
\addcontentsline{toc}{chapter}{Acknowledgements}

This thesis represents the culmination of a deeply fulfilling academic journey, and it would not have been possible without the profound support of my mentors, colleagues, and family. 

First and foremost, I gratefully acknowledge the guidance, patience, and critical feedback of my academic supervisor and the thesis defence committee at the Higher School of Economics (HSE) Master's Program in Data Science. Your rigorous standards and technical insights have fundamentally shaped my approach to research. Reflecting on my academic development, the foundational engineering principles I acquired during my Bachelor of Engineering in Mechatronics, Robotics, and Automation Engineering with a minor in Computing from the National University of Sciences and Technology (NUST) have been absolute prerequisites in grounding my deep learning models in real-world physical and kinematic boundaries.

I also extend my sincere appreciation to my co-founders and colleagues at LangMeet AI. Serving as Technical Lead while managing real-time speech translation infrastructure concurrently with my thesis workload has been a formidable challenge, but it provided an invaluable sandbox for investigating computational latency profiles, hardware-level optimization, and production bottlenecks.

On a deeply personal level, I owe an immense debt of gratitude to my family. To my parents, your shared childhood observations, continuous encouragement, and steadfast support have been my ultimate cognitive anchor. To my sister, whose empathy, kindness, and unyielding dedication to rescuing and adopting stray cats inspire me daily to notice harmony in chaotic systems. To my brother, whose background studying classical violin provided a beautiful, resonant acoustic environment that acted as the backdrop to countless late-night debugging and writing sessions. I am also deeply thankful to my aunt for her kind nature, warmth, and steady reassurance throughout my master's timeline.

Finally, I must acknowledge the quiet but essential companionship of my cat, Alec, who spent countless hours keeping watch beside my workstations while these neural networks were designed, trained, and evaluated. 

As this Master's exploration draws to a close, I look forward to expanding this foundation. The synthetic experiments and decentralized control frameworks validated herein represent an indispensable stepping stone. Translating these direct sensorimotor networks from controlled physical constraints into adaptive physical deployments on multi-robot platforms will form the explicit technical objective of my upcoming Ph.D. program.

Supplementary materials, dataset configurations, and the full implementation codebase are maintained in the open repository: \url{https://github.com/YoloPopo/neurokin}.

\clearpage

\tableofcontents
\clearpage
\listoffigures
\addcontentsline{toc}{chapter}{List of Figures}
\clearpage
\listoftables
\addcontentsline{toc}{chapter}{List of Tables}
\clearpage

% --- Main Matter ---
\pagenumbering{arabic}

% ==========================================
% CHAPTER 1: INTRODUCTION
% ==========================================
\chapter{Introduction}
\label{chap:intro}

\section{The Challenge of Embodiment Drift}
\label{sec:challenge_embodiment_drift}

Autonomous robots deployed in remote, unstructured environments---planetary rovers, disaster-response machines, deep-sea manipulators, and industrial inspection drones---must operate for extended periods without human intervention. These systems cannot rely on frequent manual recalibration, hardware replacement, or human-in-the-loop supervision. They must be self-sufficient, adaptive, and resilient in the face of inevitable physical degradation \cite{bongard2006resilient, cully2015robots}.

A fundamental requirement for such autonomy is maintaining an accurate internal representation of one's own morphology, kinematics, and dynamics. Traditional robotic systems rely on predefined kinematic models---typically specified in Unified Robot Description Format (URDF) files---that encode link lengths, joint limits, inertial properties, and collision geometries. These models enable classical control algorithms, such as Jacobian-based inverse kinematics and operational space control, to map intended high-level actions into precise joint torques or positional trajectories \cite{siciliano2009robotics}. Under ideal conditions, with a perfectly calibrated robot operating in a controlled environment, these methods achieve millimeter-level precision.

However, the assumption that a physical robot remains mathematically identical to its idealized digital model is fundamentally flawed under continuous operational deployment. Physical embodiment subjects robots to the relentless realities of thermodynamics, material fatigue, and mechanical friction. Over extended operational cycles---spanning hours, days, or months of continuous operation---the strict mathematical guarantees of forward and inverse kinematics inevitably decouple from physical reality.

This persistent, accumulating divergence between the internal computational representation of the robot and its actual physical state is formally defined in this thesis as \textit{embodiment drift} \cite{bongard2006resilient, kwiatkowski2019task}. The term captures the essential phenomenon: the robot's sense of its own body drifts away from truth as hardware degrades.

\subsection{Manifestations of Embodiment Drift}
\label{subsec:manifestations_embodiment_drift}

Embodiment drift manifests across a broad spectrum of mechanical and sensory degradations. At the micro-scale, geared transmissions develop unmodeled backlash as teeth wear; actuator drive belts lose optimal tension through creep and fatigue; thermal expansion induced by sustained high-torque operations alters the effective physical length of metallic linkages by micrometers---enough to matter in precision tasks. At the sensory layer, proprioceptive joint encoders may imperceptibly slip from their calibrated zero-points due to vibration or electrical noise; exteroceptive sensors, such as chassis-mounted RGB cameras or depth sensors, experience extrinsic calibration shifts from ambient thermal cycling or minor physical impacts \cite{rezvani2024robust}.

In more severe, sudden operational scenarios, robots may suffer macroscopic structural damage: an unexpected collision may plastically deform a physical link by several degrees; an actuator may seize entirely due to electrical failure or debris ingress; an asymmetrical payload shift---such as a grasped object suddenly changing weight distribution---may fundamentally alter the system's center of mass and dynamic response.

\subsection{Consequences for Control}
\label{subsec:consequences_control}

When significant embodiment drift occurs, the consequences for closed-loop feedback control are catastrophic. A controller operating under the assumption of a nominal kinematic tree will compute pseudo-inverse Jacobian updates that are mathematically sound relative to the model but physically erroneous relative to reality \cite{siciliano2009robotics}. The system issues corrective motor commands that inadvertently drive the end-effector away from the target, induce highly unstable resonant oscillations, or, in contact-rich manipulation tasks, exert unsafe, unbounded forces upon the environment. In legged locomotion, embodiment drift leads to gait asymmetry, reduced stability margins, and eventual falls. In humanoid robots operating near humans, the safety implications are even more severe \cite{kwiatkowski2019task}.

To maintain true autonomy and operational safety in the presence of embodiment drift, a robotic system must possess the capacity for \textit{self-modeling}---the ability to autonomously infer, update, and deploy an internal representation of its own morphology derived strictly from its own sensorimotor experience, entirely bypassing the need for external, human-in-the-loop manual recalibration \cite{chen2022full, kwiatkowski2019task}.

\begin{figure}[H]
    \centering
    \includegraphics[width=0.95\textwidth]{timeline_self_modeling.pdf}
    \caption{Three eras of robot self-modeling: symbolic morphology inference via active exploration (Bongard et al., 2006), behavioral adaptation through repertoire search (Cully et al., 2015), and contemporary visual self-modeling driven by neural rendering (NeRF and 3D Gaussian Splatting).}
    \label{fig:timeline_intro}
end{figure}

\section{The Evolution of Self-Modeling Paradigms}
\label{sec:evolution_self_modeling}

The formidable challenge of autonomous self-modeling has driven several distinct research paradigms over the past two decades. Each era brought unique insights and faced specific limitations that motivate the approach taken in this thesis. Figure~\ref{fig:timeline_intro} provides a visual summary of this progression.

\subsection{The Symbolic Era: Estimation-Exploration}
\label{subsec:symbolic_era}

The earliest formalized paradigm relied heavily on symbolic and kinematic inference. The seminal work of Bongard, Zykov, and Lipson established the first complete closed-loop self-modeling system on actual physical robots \cite{bongard2006resilient}. Their approach viewed self-modeling as a continuous cycle of hypothesis generation and active physical exploration. The robot maintained a population of competing symbolic models representing various candidate topologies---different possible configurations of its own body. To discriminate among these hypotheses, the robot systematically synthesized exploratory motor commands mathematically designed to maximize the sensory disagreement among the candidate models. Upon executing these actions, the robot collected sensor data and systematically pruned models that failed to predict observed outcomes. After approximately 16 modeling-testing cycles spanning several minutes of real-world experimentation, the robot converged on an accurate damaged model and synthesized a compensatory gait through reinforcement learning on the internal simulation.

The core innovation was an active learning loop: the fitness of a test action was determined by its ability to cause disagreement among the best candidate models, thereby maximally reducing uncertainty. While conceptually elegant, this approach suffers from the symbolic bottleneck: the robot can only discover models constructible from predefined physics engine primitives (cylinders, hinge joints). It cannot represent non-rigid deformations, soft-body dynamics, or arbitrary visual damage---limitations that become critical in unstructured environments where failure modes are unpredictable \cite{kwiatkowski2019task}. As noted by Kwiatkowski and Lipson, the search space of symbolic physics engines is discrete and non-differentiable, making optimization slow and prone to local optima \cite{kwiatkowski2019task}. Furthermore, the approach relied heavily on simple proprioceptive sensors providing limited information compared to high-bandwidth vision.

\subsection{The Behavioral Era: Avoidance Over Adaptation}
\label{subsec:behavioral_era}

Recognizing the computational intensity of explicit morphological reconstruction, a subsequent paradigm emerged focused entirely on behavioral adaptation. Rather than diagnosing the damage, these methods sought to find compensatory behaviors that worked despite the damage \cite{cully2015robots}. Cully and colleagues introduced Intelligent Trial and Error (IT\&E), which bypassed damage diagnosis entirely. Their key insight was to focus on finding what still works rather than diagnosing what went wrong.

Using the MAP-Elites quality-diversity algorithm \cite{mouret2015illuminating}, they pre-computed a behavioral repertoire of approximately 13,000 distinct gaits in simulation---a process requiring 40 million simulations over two weeks of computation. This repertoire is a high-dimensional grid where each cell corresponds to a specific behavioral descriptor (e.g., duty cycle versus leg offset) and stores the highest-performing controller found for that descriptor.

When damaged in the field, the robot treated this map as a Bayesian prior, using Gaussian Process regression to search for a behavior that performs well in the real world. By testing a few dozen distinct behaviors, the Gaussian Process updated its posterior mean and variance, rapidly identifying a compensatory gait. Cully and colleagues demonstrated recovery in less than two minutes of real-world testing across various damage scenarios, including removal of entire legs \cite{cully2015robots}.

However, IT\&E represents a strategy of avoidance, not adaptation. The robot does not repair its internal model---it discards the parts of the behavioral space that no longer function. While effective for redundant locomotion tasks where multiple gaits can achieve similar forward velocity, this approach fails for precision manipulation tasks that require exact geometric awareness \cite{kwiatkowski2019task}. If a robot needs to reach through a narrow aperture, it cannot simply ``try a different gait''---it must understand the precise Cartesian geometry of its damaged arm. Furthermore, as robotic complexity increases, the curse of dimensionality makes pre-computing dense behavioral maps prohibitively expensive: a 6-DOF arm with even modest discretization would require billions of map entries, far exceeding computational budgets.

\subsection{Task-Agnostic Learned Dynamics Models}
\label{subsec:task_agnostic_models}

The advent of deep learning enabled a third paradigm: learning dynamics models that encode the mapping from motor commands to sensory consequences. Kwiatkowski and Lipson proposed task-agnostic self-modeling, training deep forward models to predict the future state of the robot given the current state and action \cite{kwiatkowski2019task}. This represented a pivotal shift from symbolic physics to learnable physics. The same approach was demonstrated to work as a complete reinforcement learning environment to learn zero-shot tasks across morphologies in a follow-up study \cite{kwiatkowski2019zeroshot}.

While theoretically powerful, these data-driven models face a fundamental limitation: systematic bias accumulates progressively across extended planning horizons \cite{kwiatkowski2019task}. Simulated trajectories diverge increasingly from actual robot behavior, rendering the approach impractical for tasks requiring accurate prediction beyond very short time windows. This horizon-dependent degradation motivates approaches that prioritize shorter planning windows or use learned corrections rather than pure forward simulation.

\subsection{Proprioceptive Morphology Identification}
\label{subsec:proprioceptive_identification}

A parallel thread of research focuses on proprioceptive identification---restoring embodiment using internal signals alone, without the occlusions and lighting failures that often cause vision-only pipelines to break \cite{hu2024reconfigurable}. The internal sensing avoids external disturbances, offering significant deployment advantages: the approach remains robust across diverse lighting conditions and operates without calibrated cameras. However, eliminating external visual information imposes real functional costs: expressiveness is substantially reduced compared to vision-based methods \cite{farghdani2025robust}. This is not merely a theoretical gap but a genuine operating constraint that limits the complexity of tasks these systems can support compared to their multi-modal counterparts.

Meta-self-modeling restores morphology parameters between motion signature reconfigurations \cite{hu2024reconfigurable}, and IMU-centric methods show hardware-realized damage sensing with minimal sensing overhead on legged robots \cite{farghdani2025robust}. Analytical modular pipelines remain strong baselines in the case of known assembly constraints, offering interpretable, stable, and automatically generated kinematic and dynamic models \cite{dogra2022unified}. The frontier is intervention, not just identification: controlled breakage of proprioception-driven morphology demonstrates that active structural adaptation is supported by internal sensing, and that morphology itself is a form of adaptive feedback control \cite{xu2026antenna, yu2026kinematicself}.

\subsection{The Visual Era: Deep Self-Models and Neural Rendering}
\label{subsec:visual_era}

The current state-of-the-art leverages neural rendering to learn visual self-models directly from camera observations. Chen and colleagues introduced occupancy query learning as a foundational approach \cite{chen2022full}. Signed Distance Functions (SDFs) conditioned on joint angles enable differentiable collision checking directly from RGB-D sensor streams, supporting both motion planning and recovery from damage. Their multi-view hardware system achieved approximately one percent workspace accuracy and demonstrated support for planning and post-damage recovery. However, practical deployment introduced significant challenges: careful camera placement became critical for observability, precise calibration was required for geometric accuracy, and performance degraded substantially under occlusion, shadows, blur, and noise \cite{chen2022full}.

Hu, Lin, and Lipson reduced the camera system to a single RGB sensor while learning the Free-Form Kinematic Self-Model (FFKSM) from a single viewpoint \cite{hu2024teaching}. The reparameterization of camera frames in terms of the first two joints facilitated the learning process. This approach supports inverse kinematics, path planning, and post-damage recovery. However, volumetric rendering via NeRF remained computationally expensive, and single-view training introduced persistent vulnerabilities: occlusion sensitivity and perspective-dependent predictions.

More recently, Kinematic 3D Gaussian Splatting (K-3DGS) has emerged as an explicit alternative to implicit NeRFs \cite{hu2025learning}. 3DGS replaces ray marching with differentiable rasterization of explicit anisotropic Gaussians, enabling high frame rate rendering and fast optimization. By attaching explicit Gaussian primitives to a differentiable kinematic chain and leveraging rasterization-based rendering rather than ray-marching, K-3DGS promises faster inference. Each Gaussian primitive is defined by a 3D mean position, covariance matrix, opacity, and color, allowing the model to render the robot in any arbitrary pose through standard graphics pipeline operations.

Yet, as this thesis establishes, the immense representational power of these neural rendering architectures has introduced a severe computational bottleneck that completely obstructs physical deployment in real-time control loops.

\section{The Latency Paradox in Neural Rendering}
\label{sec:latency_paradox}

A diagnostic self-model is operationally valuable only if it can provide state estimates fast enough to influence the stabilizing behavior of the feedback controller \cite{hu2024teaching, hu2025learning}. If the model requires hundreds of milliseconds to generate a prediction, the information arrives too late---the robot has already executed several control cycles based on stale assumptions.

\subsection{Temporal Constraints of Robotic Control}
\label{subsec:temporal_constraints}

Modern robotic control systems operate at tightly constrained frequencies. Force control, impedance control, and dynamic balancing typically run at 1 kHz, requiring the entire perception-planning-actuation pipeline to execute within 1.0 millisecond \cite{siciliano2009robotics}. Table~\ref{tab:control_loop_latency} shows the latency requirements for different control regimes.

\begin{table}[H]
\centering
\caption{Latency requirements for robotic control loops at different frequencies.}
\label{tab:control_loop_latency}
\resizebox{\textwidth}{!}{%
\begin{tabular}{lc}
\toprule
\textbf{Control Type} & \textbf{Frequency} & \textbf{Max Latency (ms)} \\
\midrule
Force/Impedance Control & 1 kHz & 1.0 \\
Compliant Manipulation & 500 Hz & 2.0 \\
Visual Servoing & 30--100 Hz & 10.0--33.3 \\
Offline Diagnostics & <10 Hz & >100.0 \\
\bottomrule
\end{tabular}
}
\end{table}

Within the 1.0 ms window, the entire software stack must read physical sensor buffers, execute state estimation, compute target Cartesian error vectors, resolve complex inverse kinematic equations, enforce physical joint limits, and issue subsequent motor torque commands. Consequently, any self-modeling diagnostic component integrated into this inner control loop possesses an absolute maximum inference budget of a fraction of a millisecond.

\subsection{The Computational Bottleneck}
\label{subsec:computational_bottleneck}

Current state-of-the-art neural rendering architectures fundamentally violate this latency budget. Implicit NeRF-based models achieve frame rates of approximately 5 FPS, corresponding to 191.6 ms latency \cite{hu2024teaching}. Explicit 3DGS-based architectures reach approximately 37 FPS, corresponding to 27.0 ms latency \cite{hu2025learning}. Both exceed the 1.0 ms control budget by factors of 27 to 190---rendering them entirely unsuitable for real-time control integration.

The root cause is architectural. Both approaches maintain an explicit or implicit 3D volumetric representation. NeRF requires evaluating a deep network hundreds of thousands of times per image via ray-marching---a fundamentally serial process \cite{mildenhall2020nerf}. Even with acceleration structures like Instant-NGP's hash encoding \cite{muller2022instant} or Mip-NeRF's scale-aware anti-aliasing \cite{barron2021mip}, the latency remains too high for tight control loops required in robotic damage recovery. Adapting NeRFs to articulated bodies via D-NeRF approaches \cite{pumarola2021d} adds time-deformation networks, further compounding computational cost.

3DGS requires computing forward kinematics for thousands of Gaussian primitives, projecting them into 2D screen space, sorting them by depth, and alpha-compositing \cite{kerbl20233d}---a heavily sequential pipeline despite GPU acceleration. While 3DGS achieves substantial speedup over NeRF, this speed is achieved by retreating to isotropic constraints that yield geometric artifacts, dropping fidelity \cite{hu2025learning}. These operations are not merely engineering inefficiencies; they are structural consequences of demanding a full 3D intermediate representation.

\subsection{The Plasticity-Stability Dilemma in Self-Modeling}
\label{subsec:plasticity_stability}

Beyond latency, self-modeling faces the plasticity-stability dilemma, first articulated by Grossberg \cite{grossberg1987competitive}: the system must be plastic enough to learn new configurations (damage) yet stable enough to retain knowledge of undamaged components. In biological brains, this balance is achieved through synaptic pruning---the physical restructuring of the connectome wherein weak neuronal connections are eliminated while frequently used connections are strengthened. Hebbian learning (neurons that fire together, wire together) ensures that functional modules are physically distinct.

In standard, monolithic deep neural networks, sequentially learning new spatial geometries leads inevitably to catastrophic forgetting \cite{mccloskey1989catastrophic, french1999catastrophic}. If an implicit NeRF, an explicit 3DGS point cloud, or a direct convolutional decoder is naively fine-tuned exclusively on a stream of recent visual observations depicting a damaged joint, the globally shared gradient updates will rapidly destroy the network's ability to accurately render the undamaged base of the robot. The network overfits to the new, damaged data distribution at the catastrophic expense of prior structural knowledge \cite{kirkpatrick2017overcoming}.

To combat this sequential learning degradation, the literature has developed extensive mitigation frameworks. Regularization-based methodologies, such as Elastic Weight Consolidation (EWC) \cite{kirkpatrick2017overcoming}, calculate the Fisher information matrix to identify and mathematically protect specific neural weights crucial for maintaining prior geometric knowledge. Architectural isolation methods rely on dynamic routing protocols, mixture-of-experts ensembles \cite{mocanu2018scalable}, or isolating continuous learning within highly sparse, trainable subnetworks identified via the Lottery Ticket Hypothesis \cite{frankle2018lottery}. However, these methods remain computationally expensive for real-time adaptation on edge hardware, and they treat the network as a monolith with varying ``stiffness'' rather than enabling true structural modularity.

\section{Digital Twins and the Simulation-to-Reality Gap}
\label{sec:digital_twins}

The endeavor to maintain an accurate, real-time internal self-model is conceptually identical to the broader industrial engineering goal of generating and sustaining a robotic Digital Twin \cite{jiang2024digitaltwin, mari2024digitaltwin}. In contemporary academic literature, a digital twin is defined not merely as a static 3D visual replica or a passive CAD model, but as a bi-directional, dynamically synchronized computational surrogate of a physical asset.

The idealized robotic digital twin operates in continuous tandem with the physical hardware. It receives high-frequency sensor telemetry (joint angles, motor torques, camera feeds) from the physical robot in the field. Crucially, it dynamically updates its internal state matrices---its collision geometry, mass distribution, and friction coefficients---to perfectly mirror the physical degradation and embodiment drift experienced by the hardware. This dynamic synchronization allows the twin to serve as a perfectly safe, highly accelerated sandbox for forward-simulating potential control policies before executing them in physical reality \cite{sun2026digitaltwin, murugesan2024digitaltwin}.

\begin{figure}[H]
    \centering
    \includegraphics[width=0.95\textwidth]{sim2real_digital_twin.pdf}
    \caption{Digital twin synchronization loop bridging sensing, simulation, and control.}
    \label{fig:digital_twin_loop}
\end{figure}

However, constructing visual digital twins specifically intended for autonomous, closed-loop robotic control introduces stringent latency constraints that differ significantly from those found in static industrial monitoring. An active robotic digital twin must be instantaneously queryable. If a physical manipulator operating in a highly cluttered environment sustains an unexpected impact, the digital twin must update its geometric occupancy map immediately. If the twin's internal update latency exceeds the temporal dynamics of the physical interaction, the robot will execute evasive commands based on stale, pre-impact geometry, precipitating further catastrophic collisions.

Furthermore, the practical deployment of neural self-models and dynamic digital twins is universally haunted by the Simulation-to-Reality (Sim-to-Real) gap \cite{tan2018simtoreal, aljalbout2026realitygap}. The majority of high-capacity neural representations are trained within pristine simulation environments, such as MuJoCo, Isaac Gym, or, as utilized in this thesis, PyBullet. While computationally efficient, these physics simulators rely on idealized mathematical approximations. They frequently fail to accurately capture complex Coulomb friction models, non-linear actuator stiction, subtle link compliance under heavy payloads, and the incredibly complex photometric realities of physical environments, including ambient lighting fluctuations, harsh specular reflections on metallic links, and pervasive sensor noise.

When a neural self-model trained exclusively in simulation is deployed onto physical hardware, it frequently collapses. The minor photometric discrepancies between the simulated training data and the physical camera feed severely corrupt the network's latent space, resulting in completely erroneous morphological predictions \cite{rezvani2024robust}. Recent literature attempts to bridge this critical Sim-to-Real gap through massive Domain Randomization, where researchers randomly perturb lighting vectors, background textures, camera positions, and physical mass properties during simulation training \cite{tan2018simtoreal}. Advanced methodologies also explore real-time Residual Dynamics Learning, wherein the neural network is trained to predict only the small, non-linear residual discrepancy between the idealized physics simulator and the actual physical observation \cite{qiu2024swim2real, lim2025splat2real}.

Within the rigorous scope of this thesis, all experiments are deliberately bounded within the synthetic PyBullet simulation environment. This is not an academic oversight but a highly controlled, necessary methodological choice designed specifically to isolate architectural latency. By operating within the controlled vacuum of a simulator, this thesis effectively isolates the theoretical mathematical efficiency of the proposed NeuroKin architecture, proving that the sub-millisecond bound can be breached before introducing the complexities of physical noise.

\section{Problem Statement and Core Hypothesis}
\label{sec:problem_hypothesis_intro}

The core scientific problem addressed by this thesis is the resolution of the latency paradox in visual robotic self-modeling, and the subsequent leveraging of this high-speed diagnostic capability to establish the computational mechanisms for multi-agent swarm coordination under compounding mechanical failures.

\subsection{Formal Mathematical Problem Statement}
\label{subsec:formal_mathematical_problem_statement_intro}

We formally define the true physical morphology of the robot at time $t$ as the configuration manifold $\mathcal{M}(t)$. The robotic system maintains an internal computational kinematic representation $\hat{\mathcal{M}}(t)$. Embodiment drift is defined as the non-stationary, cumulative divergence between the physical reality and the digital expectation:
\begin{equation}
D(t) = \|\mathcal{M}(t) - \hat{\mathcal{M}}(t)\|
\end{equation}

When an unmodeled kinematic fault $\Delta \mathbf{q}_{\text{fault}}$ occurs, the proprioceptive encoder estimate $\mathbf{q}_{\text{enc}}$ deviates from the true physical state $\mathbf{q}_{\text{true}}$ such that $\mathbf{q}_{\text{true}} = \mathbf{q}_{\text{enc}} + \Delta \mathbf{q}_{\text{fault}}$. This breaks traditional closed-loop control. 

Let the sensory observation space be $\mathcal{O}$ (in this case, binary visual silhouettes) and the joint configuration space be $\mathcal{Q}$. The objective of the visual self-model is to learn a direct sensorimotor mapping $f_\theta: \mathcal{Q} \rightarrow \mathcal{O}$ that minimizes the reconstruction loss $\mathcal{L}(f_\theta(\mathbf{q}), \mathbf{o})$ to recover the true Cartesian end-effector position $\mathbf{x}_{\text{true}}$. 

Crucially, this mapping is strictly bound by control-loop temporal dynamics. If the control system operates at frequency $f_{ctrl}$, dictating a strict maximum latency constraint $T_{max} = 1/f_{ctrl}$, and $t_{\text{infer}}$ is the inference time of the self-model, the system is physically deployable for closed-loop stabilization if and only if:
\begin{equation}
t_{\text{infer}} \le T_{max}
\end{equation}
For a standard $1$~kHz feedback loop, this necessitates $t_{\text{infer}} \le 1.0$~ms. 

\subsection{The Central Hypothesis}
\label{subsec:central_hypothesis_intro}

This thesis investigates a central hypothesis: \textit{explicit 3D volumetric reconstruction is not a prerequisite for control-oriented self-modeling}. If the downstream task requires only silhouette-level awareness, collision checking, or end-effector coordinate estimation from known camera viewpoints, maintaining a dense, queryable 3D scene representation is computationally wasteful. The robot does not need to know the precise 3D occupancy of every voxel in its volume; it needs to know where its end-effector is in Cartesian space and whether its links will collide with obstacles.

We propose a paradigm shift toward \textit{direct sensorimotor decoding}: a lightweight fully convolutional network that maps joint angles directly to visual silhouettes and spatial coordinates, completely bypassing the intermediate 3D reconstruction step. We hypothesize that such an architecture can achieve sub-millisecond inference latency---well within the 1 kHz control budget---while maintaining reconstruction fidelity comparable to or exceeding volumetric baselines, precisely because dense pixel-wise supervision is more efficient than sparse ray sampling.

\subsection{Empirical Scope, Constraints, and Hardware Realities}
\label{subsec:empirical_scope_intro}

The empirical scope of this research is deliberately bounded within a highly controlled simulation environment to isolate core architectural latency metrics from unmodeled photometric noise. All data generation, neural network training, closed-loop controller validation, and multi-agent coordination are executed using a synthetic PyBullet simulation of a 4-DOF serial manipulator. The robot is rendered against a uniform background; joint angles are recorded without sensor noise; camera calibration is perfect. These constraints are not limitations but methodological choices: they allow us to measure pure architectural latency without confounding variables.

This boundary condition is a deliberate necessity. Implementation on actual physical hardware is not possible within the scope of this Master's thesis due to severe, practical resource limitations. The rigorous validation of dynamic fault adaptation requires dedicated, damageable physical manipulators where structural faults can be safely and repeatedly injected. Furthermore, acquiring ground-truth millimetric spatial coordinate data to validate the visual network's predictions requires expensive, high-speed motion-capture rigs (e.g., Vicon systems). Given the strict temporal and hardware constraints of the degree program, physical deployment is unattainable at this stage. Consequently, the transition from simulated architectural validation to physical hardware deployment---including addressing the complex photometric noise of the Sim-to-Real gap---is formally deferred to the author's upcoming Ph.D. research.

All experiments use a dataset of 2,000 samples (1,600 training, 400 validation) generated via Lorenz attractor trajectories, ensuring smooth, continuous, and diverse workspace coverage. Hardware is an NVIDIA Tesla T4 GPU with PyTorch 2.0 and CUDA 11.8.

\subsection{The Temporal Stigmergy Framework for Swarm Coordination}
\label{subsec:temporal_stigmergy_framework_intro}

Crucially, the thesis title---``Toward Swarm Coordination''---is framed with scientific restraint. This thesis does not claim concurrent spatial multi-agent deployment. Deploying multiple physical robots simultaneously in a shared workspace introduces immense additional complexities: radio-frequency bandwidth constraints, physical collision avoidance algorithms, decentralized simultaneous localization and mapping, and distributed task allocation \cite{matsuki2024gaussian, keetha2024splattam}. These are separate research problems beyond this thesis's scope.

Instead, we evaluate multi-agent coordination via \textit{temporal stigmergy}. Stigmergy is a biological mechanism of indirect coordination observed in social insects such as termites and ants \cite{bongard2006resilient}. Individual agents do not communicate directly; instead, they alter the environment in ways that stimulate specific responses from subsequent agents. A pheromone trail deposited by one ant triggers following behavior in others. A pellet of mud placed by one termite becomes a stimulus for another to add more mud at that location.

In this thesis, the shared ``environment'' is algorithmically abstracted into a persistent digital memory ledger---the \texttt{factory\_state} variable. Agents do not broadcast distress signals or negotiate global plans. Instead, when an agent experiences repeated mechanical failures under compounding faults, it updates the \texttt{consecutive\_failures} integer in the shared ledger. Subsequent agents, upon initiation, read this trace. If the trace indicates a degraded operational history (two consecutive failures), the new agent autonomously shifts into a localized ``recovery'' mode: it abandons its randomized spatial target, moves to a conservative centralized coordinate, and doubles its Cartesian error tolerance.

This state-dependent, history-aware behavioral adaptation is the absolute core logic of decentralized swarm coordination. Proving that sub-millisecond individual self-diagnosis enables such temporal coordination validates the computational prerequisite for eventual physical swarms.

\section{Research Questions}
\label{sec:research_questions_intro}

Four research questions guide this investigation, progressing from architectural efficiency through control integration and fault tolerance to collective adaptation.

\subsection{RQ1: Architectural Efficiency}
\label{subsec:RQ1_intro}

To what extent can the computational latency of visual robot self-modeling be reduced by replacing 3D implicit and explicit volumetric rendering paradigms with 2D direct sensorimotor decoding, and what is the resulting impact on morphological reconstruction fidelity measured in PSNR? This question probes the fundamental architectural trade-off between representational complexity and inference speed, seeking to establish whether direct decoding can simultaneously improve both metrics.

\subsection{RQ2: Control Integration}
\label{subsec:RQ2_intro}

Can a visual self-model operating purely on direct convolutional decoding provide spatial coordinate estimates with sufficient accuracy and sub-millisecond latency to actively stabilize a Jacobian-based inverse kinematics control loop? This question moves beyond static rendering metrics to dynamic closed-loop performance, asking whether the network's predictions are smooth and accurate enough to serve as the primary state estimate for feedback control.

\subsection{RQ3: Fault Tolerance}
\label{subsec:RQ3_intro}

How does the integration of sub-millisecond visual self-estimates alter end-effector trajectory and error convergence compared to a purely proprioceptive baseline under induced mechanical faults such as severe encoder slip? This question evaluates the system's resilience under the very conditions that motivate self-modeling in the first place: when physical damage corrupts proprioceptive readings, can visual estimates override them?

\subsection{RQ4: Collective Adaptation}
\label{subsec:RQ4_intro}

Can the individual computational efficiency of a sub-millisecond self-model enable decentralized, state-dependent adaptation across a sequential pipeline of agents, thereby validating the temporal stigmergy mechanics required for future multi-agent swarm coordination? This question scales the investigation from single-agent survival to collective endurance.

\section{Methodological Contributions}
\label{sec:methodological_contributions_intro}

This thesis makes four primary contributions to adaptive robotics, neural representation, and decentralized multi-agent systems.

\subsection{Contribution 1: Formal Benchmarking of the Latency Paradox}
\label{subsec:contrib_benchmarking_intro}

The first contribution is the formal benchmarking of the latency paradox in neural rendering. Through rigorous ground-up reproduction and quantitative evaluation of both implicit (FFKSM/NeRF) and explicit (K-3DGS) neural rendering self-models on a standardized dataset, this research formally exposes their control incompatibility. The exact latency metrics---191.6 ms for FFKSM, 27.0 ms for K-3DGS---establish the mathematical limits that necessitate a paradigm shift.

\subsection{Contribution 2: The NeuroKin Direct Sensorimotor Decoder}
\label{subsec:contrib_neurokin_intro}

The second and most significant contribution is the introduction and exhaustive evaluation of the NeuroKin architecture. NeuroKin is a lightweight, fully convolutional direct sensorimotor decoder that completely bypasses volumetric reconstruction, mapping joint angles directly to visual silhouettes and spatial coordinates. Through extensive empirical validation, NeuroKin achieves 21.88 dB PSNR at 7,400 FPS (0.135 ms latency)---a 1,418$\times$ speedup over implicit NeRF baselines with superior fidelity. This shatters the sub-millisecond inference bottleneck, rendering high-fidelity visual self-modeling compatible with 1 kHz control loops for the first time. An auxiliary multi-task variant, ResNeuroKin-D, recovers geometric depth at 2,500 FPS (0.400 ms latency).

\subsection{Contribution 3: Visual Closed-Loop Damage Adaptation}
\label{subsec:contrib_closed_loop_intro}

The third contribution lies in visual closed-loop control and dynamic damage adaptation. This research successfully synthesizes NeuroKin's sub-millisecond spatial estimates---derived via a synchronized stereoscopic dual-camera decoding pipeline---with a Jacobian-based inverse-kinematics feedback loop. Through empirical demonstrations of dynamic error recovery following simulated structural damage, the thesis proves that high-speed visual estimates can override corrupted proprioceptive data, allowing the controller to re-stabilize the end-effector trajectory post-fault.

\subsection{Contribution 4: Temporal Stigmergy for Sequential Swarm Coordination}
\label{subsec:contrib_stigmergy_intro}

The fourth contribution is the validation of stigmergic swarm coordination under compounding mechanical faults. The thesis executes a 100-stage continuous sequential swarm simulation under escalating fault injection probabilities. The standard inverse-kinematic baseline collapses to 13\% post-fault success, while the NeuroKin-enabled sequence maintains 88\% success by reading the shared stigmergic ledger and triggering decentralized recovery modes. This supports the claim that ultra-fast individual self-diagnosis is a practical prerequisite for more capable swarm coordination.

\section{Thesis Organization}
\label{sec:thesis_organization_intro}

The remainder of this thesis is structured to systematically address the research questions. Chapter 2 provides an exhaustive critical synthesis of relevant literature, tracing the mathematical evolution of robot self-modeling, analyzing the constraints of NeRF and 3DGS, and establishing the theoretical framework of temporal stigmergy for swarm robotics.

Chapter 3 details the empirical simulation methodology and latency bottleneck, outlining the PyBullet data generation pipeline, the application of Lorenz attractors for chaotic workspace exploration, and providing rigorous quantitative breakdowns of the FFKSM and K-3DGS baseline reproductions.

Chapter 4 formally introduces the core technical innovation---the NeuroKin architecture---detailing its fully convolutional network topology, stereoscopic spatial recovery mechanisms, and the definitive performance evaluation proving the 7,400 FPS inference breakthrough.

Chapter 5 moves from static rendering metrics to dynamic control, outlining the integration of NeuroKin into an inverse-kinematics loop and providing step-by-step analysis of end-effector convergence under both nominal and fault-injected conditions.

Chapter 6 scales validation to the collective level, unpacking the 100-stage operational swarm chain, explaining the mechanics of the inter-agent communication ledger, and analyzing the statistical success rates of the temporal stigmergy pipeline against cascading baseline failures.

Chapter 7 critically evaluates the overarching hypothesis, analyzing the Pareto trade-offs of sacrificing novel-view synthesis for extreme inference speed, defending the stigmergic approach to swarm coordination, and discussing critical threats to validity---particularly the simulation-to-reality gap and the abstraction of mechanical degradation.

Finally, Chapter 8 synthesizes the research findings, directly answers the formal research questions, and maps a concrete technological roadmap for transitioning these validated sequential capabilities into physically embodied, spatially concurrent swarm deployments.
\clearpage

% ==========================================
% CHAPTER 3: BASELINE IMPLEMENTATIONS AND THE LATENCY BOTTLENECK
% ==========================================
\chapter{Baseline Implementations and the Latency Bottleneck}
\label{chap:baseline_latency_bottleneck_ch3}

To rigorously characterize the latency paradox that motivates this thesis, we implement and evaluate two state-of-the-art visual self-modeling architectures: the Free-Form Kinematic Self-Model (FFKSM) based on implicit neural radiance fields \cite{hu2024teaching}, and Kinematic 3D Gaussian Splatting (K-3DGS) based on explicit geometric primitives \cite{hu2025learning}. Both are reproduced from published specifications using an identical dataset, hardware environment, and evaluation protocol. This controlled comparison establishes the quantitative baseline against which our proposed NeuroKin architecture is measured.

\section{Synthetic Data Generation via Lorenz Attractors}

All experiments use a dataset of 2,000 samples generated from a PyBullet simulation of a 4-DOF serial manipulator. This section details the complete data generation pipeline, including simulation configuration, trajectory planning, and preprocessing procedures.

\subsection{Simulation Environment Configuration}

The PyBullet physics simulator provides a controlled environment for generating ground-truth robot observations. We configure the simulation with the following parameters:

\begin{itemize}
    \item \textbf{Robot Model}: A 4-degree-of-freedom serial manipulator with revolute joints. The base is fixed at the origin with joint axes arranged in an anthropomorphic configuration: joint 1 rotates about the vertical axis (yaw), joints 2 and 3 rotate about horizontal axes (pitch), and joint 4 provides wrist rotation.
    \item \textbf{Camera Setup}: A fixed RGB camera positioned at coordinates $[1.0, 0.0, 0.0]$ with focal length 130.2545. The camera looks toward the origin, producing $100 \times 100$ pixel images with a field of view of 42 degrees.
    \item \textbf{Rendering}: The simulation runs in DIRECT mode (no graphical display) for maximum speed. Each step renders the robot using PyBullet's built-in tiny renderer with shadows disabled.
    \item \textbf{Physics}: Gravity is set to $-9.8 \, \text{m/s}^2$ in the z-direction. A ground plane is placed at $z = -0.109$ m. Each joint is position-controlled with a maximum force of 100 N.
\end{itemize}

The robot's kinematic structure follows the Denavit-Hartenberg (DH) convention. The DH parameters are detailed in Table 3.1.

The forward kinematics transformation from base to end-effector is given by the product of individual link transformations:

\begin{equation}
T_{0}^{4} = T_{0}^{1}(\theta_1) T_{1}^{2}(\theta_2) T_{2}^{3}(\theta_3) T_{3}^{4}(\theta_4)
\end{equation}

where each $T_{i-1}^{i}$ is computed using the standard DH convention:

\begin{equation}
T_{i-1}^{i} = 
\begin{bmatrix}
\cos\theta_i & -\sin\theta_i & 0 & a_{i-1} \\
\sin\theta_i \cos\alpha_{i-1} & \cos\theta_i \cos\alpha_{i-1} & -\sin\alpha_{i-1} & -d_i \sin\alpha_{i-1} \\
\sin\theta_i \sin\alpha_{i-1} & \cos\theta_i \sin\alpha_{i-1} & \cos\alpha_{i-1} & d_i \cos\alpha_{i-1} \\
0 & 0 & 0 & 1
\end{bmatrix}
\end{equation}

\subsection{Lorenz Attractor Trajectory Generation}

Rather than relying on random motor babbling---which produces jerky, discontinuous motions that poorly sample the workspace---we employ trajectories generated by the Lorenz dynamical system. The Lorenz attractor is a system of three coupled ordinary differential equations (ODEs) that exhibits deterministic chaos, producing trajectories that are smooth, non-repeating, and ergodic within a bounded region of state space.

The Lorenz system is defined as:

\begin{align}
\frac{dx}{dt} &= \sigma(y - x) \\
\frac{dy}{dt} &= x(\rho - z) - y \\
\frac{dz}{dt} &= xy - \beta z
\end{align}

The parameter values determine the dynamical behavior. We use the classical values $\sigma = 10$, $\rho = 28$, $\beta = 8/3$, which place the system in the chaotic regime. At these parameters, the Lyapunov exponent is positive, ensuring sensitive dependence on initial conditions and exponential divergence of nearby trajectories. The system has two unstable fixed points at $(\pm\sqrt{\beta(\rho-1)}, \pm\sqrt{\beta(\rho-1)}, \rho-1)$ and a stable fixed point at the origin (which is only stable for $\rho < 1$). For $\rho = 28$, the origin is a saddle point, and trajectories are attracted to a strange attractor---a fractal set of measure zero with non-integer Hausdorff dimension.

To numerically integrate the Lorenz system, we employ the fourth-order Runge-Kutta (RK4) method. For a general ODE $\frac{d\mathbf{x}}{dt} = \mathbf{f}(\mathbf{x}, t)$, the RK4 update from state $\mathbf{x}_n$ at time $t_n$ to $\mathbf{x}_{n+1}$ at time $t_{n+1} = t_n + \Delta t$ is:

\begin{align}
\mathbf{k}_1 &= \mathbf{f}(\mathbf{x}_n, t_n) \\
\mathbf{k}_2 &= \mathbf{f}(\mathbf{x}_n + \frac{\Delta t}{2}\mathbf{k}_1, t_n + \frac{\Delta t}{2}) \\
\mathbf{k}_3 &= \mathbf{f}(\mathbf{x}_n + \frac{\Delta t}{2}\mathbf{k}_2, t_n + \frac{\Delta t}{2}) \\
\mathbf{k}_4 &= \mathbf{f}(\mathbf{x}_n + \Delta t \mathbf{k}_3, t_n + \Delta t) \\
\mathbf{x}_{n+1} &= \mathbf{x}_n + \frac{\Delta t}{6}(\mathbf{k}_1 + 2\mathbf{k}_2 + 2\mathbf{k}_3 + \mathbf{k}_4)
\end{align}

RK4 is a fourth-order method, meaning the local truncation error is $O(\Delta t^5)$ and the global error is $O(\Delta t^4)$. This provides excellent stability and accuracy for the Lorenz system, which is stiff in certain regions of state space.

We set the integration time step to $\Delta t = 0.01$ seconds. This value was chosen through systematic experimentation: larger steps ($\Delta t \geq 0.02$) caused numerical instability and trajectory divergence, while smaller steps ($\Delta t \leq 0.005$) increased computational cost without measurable improvement in trajectory quality.

After each RK4 step, the raw Lorenz state values are scaled using the hyperbolic tangent function:

\begin{equation}
\text{action} = \tanh(\alpha \cdot \mathbf{x})
\end{equation}

where $\alpha = 0.025$ is a scaling factor. The $\tanh$ function maps the unbounded Lorenz state to the bounded interval $[-1, 1]$, which corresponds to the normalized joint angle range. This scaling preserves the sign and relative magnitude of variations while ensuring the command stays within valid limits.

Two independent Lorenz systems drive the four robot joints. Specifically:

\begin{itemize}
    \item Joint 1 follows $x$ component of trajectory 1
    \item Joint 2 follows $y$ component of trajectory 1
    \item Joint 3 follows $x$ component of trajectory 2
    \item Joint 4 follows $z$ component of trajectory 2
\end{itemize}

Using two independent attractors (with slightly different time steps: $\Delta t_1 = 0.01$, $\Delta t_2 = 0.012$) ensures that the four joints are not perfectly correlated, producing a richer diversity of poses.

\subsection{Data Collection and Preprocessing}

For each of the 2,000 samples, the data collection procedure follows:

\begin{enumerate}
    \item The robot is reset to a neutral configuration (all joints at 0 degrees).
    \item A target joint angle vector $\mathbf{q} \in [-90^\circ, +90^\circ]^4$ is computed from the Lorenz trajectories.
    \item Position controllers move each joint to the target angle over 50 simulation steps (equivalent to 200 ms of physical time), allowing the robot to settle.
    \item An RGB image is captured from the fixed camera.
    \item Joint angles are recorded from the simulation state.
\end{enumerate}

The RGB images are preprocessed to produce binary silhouette masks:

\begin{enumerate}
    \item Convert RGB to grayscale: $I_{\text{gray}} = 0.299R + 0.587G + 0.114B$
    \item Apply threshold: $I_{\text{mask}}(u,v) = \begin{cases} 1 & \text{if } I_{\text{gray}}(u,v) < 240 \\ 0 & \text{otherwise} \end{cases}$
    \item The threshold value 240 was empirically determined to optimally separate robot pixels (typically lighter due to white/gray coloring) from background (darker).
\end{enumerate}

The dataset comprises 2,000 samples split into 1,600 training samples (80\%) and 400 validation samples (20\%). The split is performed randomly with a fixed seed (42) for reproducibility.

\subsection{Dataset Statistics and Quality Analysis}

A parameter configuration matrix test evaluates the training data properties. Statistical analysis confirms good workspace coverage:

\begin{itemize}
    \item Joint angle means: $-0.23^\circ$ to $12.45^\circ$ (near zero, as expected for a system centered at neutral)
    \item Joint angle standard deviations: approximately $30^\circ$ (indicating broad exploration)
    \item Robot pixel occupancy: average 14.7\%, range 8.9\%--23.4\%
    \item Image dimensions: $100 \times 100$ pixels (10,000 total pixels per image)
\end{itemize}

The low average occupancy (14.7\%) creates severe class imbalance---85\% of pixels are background (0) and only 15\% are robot (1). This imbalance creates a pathological convergence trap for naive volumetric rendering, as the network can achieve low loss by simply predicting all pixels as 0. This phenomenon, which we term \textit{black-screen convergence}, is detailed in Section 3.2.3.

\section{Free-Form Kinematic Self-Model (FFKSM): Implicit Neural Rendering}

The FFKSM architecture, introduced by Hu, Lin, and Lipson \cite{hu2024teaching}, adapts Neural Radiance Fields (NeRF) \cite{mildenhall2020nerf} to articulated robot self-modeling. Unlike traditional NeRF, which assumes static scenes, FFKSM conditions the radiance field on joint angles, enabling dynamic rendering as the robot moves.

\subsection{Core Mathematical Framework}

\subsubsection{Neural Radiance Fields for Static Scenes}

NeRF represents a static 3D scene as a continuous function $F_\theta: (\mathbf{x}, \mathbf{d}) \rightarrow (\mathbf{c}, \sigma)$ parameterized by a multi-layer perceptron (MLP) with weights $\theta$. The function maps a 3D spatial position $\mathbf{x} \in \mathbb{R}^3$ and a viewing direction $\mathbf{d} \in \mathbb{R}^3$ (normalized to unit length) to an emitted color $\mathbf{c} \in \mathbb{R}^3$ (RGB values in $[0,1]^3$) and a volume density $\sigma \in \mathbb{R}^+$.

To render a pixel, NeRF approximates the volume rendering integral along a camera ray $\mathbf{r}(t) = \mathbf{o} + t\mathbf{d}$, where $\mathbf{o}$ is the camera origin and $t \in [t_n, t_f]$ ranges from near plane to far plane. The expected color $C(\mathbf{r})$ is:

\begin{equation}
C(\mathbf{r}) = \int_{t_n}^{t_f} T(t) \sigma(\mathbf{r}(t)) \mathbf{c}(\mathbf{r}(t), \mathbf{d}) \, dt
\end{equation}

where $T(t) = \exp\left(-\int_{t_n}^{t} \sigma(\mathbf{r}(s)) \, ds\right)$ is the accumulated transmittance. In practice, this continuous integral is approximated via numerical quadrature using stratified sampling. Along each ray, $N$ points $\{t_i\}_{i=1}^N$ are sampled, and the discrete approximation is:

\begin{equation}
\hat{C}(\mathbf{r}) = \sum_{i=1}^{N} T_i (1 - \exp(-\sigma_i \delta_i)) \mathbf{c}_i
\end{equation}

where $\delta_i = t_{i+1} - t_i$ is the distance between adjacent samples, and $T_i = \exp\left(-\sum_{j=1}^{i-1} \sigma_j \delta_j\right)$.

For a $100 \times 100$ image with $N = 64$ samples per ray, this requires $100 \times 100 \times 64 = 640,000$ network evaluations per frame. This computational burden is the fundamental bottleneck that prevents real-time control integration.

\subsubsection{Positional Encoding}

Standard MLPs tend to learn low-frequency structure more easily than sharp high-frequency detail \cite{rahaman2019spectral}. To capture high-frequency geometric details (e.g., sharp edges of robot links), NeRF employs positional encoding that maps input coordinates to a higher-dimensional space using sinusoidal functions:

\begin{equation}
\gamma(p) = [\sin(2^0 \pi p), \cos(2^0 \pi p), \ldots, \sin(2^{L-1} \pi p), \cos(2^{L-1} \pi p)]
\end{equation}

For 3D coordinates, this transforms $\mathbf{x} \in \mathbb{R}^3$ to a $3 \times (1 + 2L)$-dimensional vector. Following the original NeRF, we use $L = 5$ frequencies, yielding 33-dimensional encodings.

\subsubsection{FFKSM Adaptation for Articulated Robots}

FFKSM extends NeRF to articulated robots through two key innovations: the \textit{virtual frame transformation} and the \textit{split encoder architecture}.

\textbf{Virtual Frame Transformation:} The virtual frame transforms 3D query points using the first two joint angles ($\theta_0, \theta_1$) before network processing:

\begin{equation}
\mathbf{x}_{\text{virtual}} = R_z(-\theta_0) R_y(-\theta_1) \mathbf{x}_{\text{world}}
\end{equation}

where $R_z$ and $R_y$ are rotation matrices:

\begin{equation}
R_z(\alpha) = \begin{bmatrix}
\cos\alpha & -\sin\alpha & 0 \\
\sin\alpha & \cos\alpha & 0 \\
0 & 0 & 1
\end{bmatrix}, \quad
R_y(\beta) = \begin{bmatrix}
\cos\beta & 0 & \sin\beta \\
0 & 1 & 0 \\
-\sin\beta & 0 & \cos\beta
\end{bmatrix}
\end{equation}

This transformation reduces the network's learning burden by aligning coordinates with the robot's base orientation, effectively factoring out base rotation from the end-effector position learning problem. Without this prior, the network would need to learn circular symmetry---a challenging task for MLPs.

\textbf{Split Encoder Architecture:} FFKSM separates coordinate processing from kinematic processing:

\begin{enumerate}
    \item \textbf{Coordinate Encoder}: Takes virtual-frame 3D points via positional encoding (5 frequencies $\rightarrow$ 33 dimensions) followed by a two-layer MLP (128 hidden units each).
    \item \textbf{Kinematic Encoder}: Processes the remaining two joint angles ($\theta_2, \theta_3$) via separate positional encoding (2 $\times$ 5 frequencies $\rightarrow$ 2 $\times$ 11 = 22 dimensions) and an identical two-layer MLP (128 hidden units).
    \item \textbf{Predictive Module}: Concatenates the 128-dimensional features from both encoders (256 total) and passes through four linear layers (128 $\rightarrow$ 128 $\rightarrow$ 32 $\rightarrow$ 2) to predict density $\sigma$ and visibility $v$.
\end{enumerate}

The split architecture is critical: a single encoder processing $(\mathbf{x}, \theta)$ jointly would need to relearn the kinematic transformation for every 3D point. The split design allows the kinematic encoder to learn the mapping $\theta \rightarrow \text{link transformation}$ once, then apply it to all query points, dramatically improving sample efficiency.

\subsection{Implementation Details and Reproduction Challenges}

Our independent implementation of FFKSM from scratch uncovered three critical challenges not detailed in the original publication. Each challenge required significant debugging and architectural modifications.

\subsubsection{Challenge 1: Black-Screen Convergence}

\textbf{Symptom:} Initial training exhibited a pathological failure mode: the network converged to predicting uniformly black images (all zeros) despite achieving apparently low loss values. The loss decreased smoothly, suggesting successful optimization, but qualitative inspection revealed complete failure.

\textbf{Root Cause Analysis:} Robots occupy only $\approx$15\% of image pixels, creating severe class imbalance (85\% background, 15\% robot). The network discovered a local minimum where predicting the majority class (black) achieved mean squared error of $\sim$0.85, yet produced no useful representation. This local minimum is extremely attractive because the gradient magnitude for background pixels is near zero, and the gradient from robot pixels is sparse and weak. The optimizer follows the path of least resistance, settling into the trivial solution. This failure persisted for the first 1,000 iterations, rendering standard training completely ineffective.

\textbf{Attempted Solutions:} Loss modification experiments were systematically performed. Applying conservative $2\times$ and $3\times$ scaling factors to foreground variables proved insufficient to shift the gradients out of the background attractor. Elevating the foreground multiplier to $5\times$ yielded initial trace improvements but introduced severe numerical boundary instability. When scaling targets were pushed to $10\times$ and $20\times$, the optimizer destabilized immediately, producing volatile oscillations and rapid network mode collapse. Focal loss implementations introduced hyperparameter sweeps that proved highly sensitive, failing to break the all-zero local minimum. Continuous Dice loss metrics caused fractional boundaries to fluctuate wildly and generated catastrophic gradient explosions.

\textbf{Solution: Center-Cropping Curriculum Learning}

We solved the black-screen convergence through a structured curriculum learning approach. The key insight is to eliminate the class imbalance during critical early learning by focusing exclusively on image regions where robot occupancy is high.

For the first $N_{\text{crop}} = 500$ iterations, training focuses on the central $50 \times 50$ region where robot occupancy exceeds 40\%. After this period, training expands to the full $100 \times 100$ image.

The curriculum can be formalized as a time-varying loss mask:

\begin{equation}
\mathcal{L}_t = \begin{cases}
\frac{1}{50^2} \sum_{(i,j) \in \text{center}} w(I_{ij}^{\text{GT}}) \|I_{ij}^{\text{pred}} - I_{ij}^{\text{GT}}\|^2, & t < 500 \\
\frac{1}{100^2} \sum_{i=1}^{100}\sum_{j=1}^{100} w(I_{ij}^{\text{GT}}) \|I_{ij}^{\text{pred}} - I_{ij}^{\text{GT}}\|^2, & t \geq 500
\end{cases}
\end{equation}

where $w(I_{ij}^{\text{GT}}) = 5$ for robot pixels ($I_{ij}^{\text{GT}} = 1$) and $1$ for background pixels ($I_{ij}^{\text{GT}} = 0$).

The center crop dimensions (50$\times$50) were chosen because at this crop size, robot occupancy consistently exceeds 40\% across all training samples, the crop is large enough to capture the entire robot's articulation range, and it eliminates the class imbalance that causes black-screen convergence. The transition point (500 iterations) was empirically determined: earlier transitions (<300 iterations) reintroduced the black-screen problem, while later transitions (>800 iterations) slowed overall convergence without benefit.

\subsubsection{Challenge 2: Kinematic Blindness}

\textbf{Symptom:} After apparent convergence at iteration 2,000, all metrics indicated success: PSNR improved from 6.94 dB to 17.35 dB, validation loss decreased smoothly from 0.24 to 0.023, and training curves appeared stable. However, qualitative validation revealed a subtle but devastating bug: all poses with identical base orientation ($\theta_0, \theta_1$) produced identical predictions regardless of end-effector configuration ($\theta_2, \theta_3$). The network exhibited "kinematic blindness" to distal joints, essentially modeling only the robot's base as a rigid cylinder.

\textbf{Root Cause Analysis:} Gradient flow analysis using PyTorch's autograd hooks uncovered the root cause: a tensor slicing bug where only $(\theta_0, \theta_1)$ reached the kinematic encoder.

The erroneous code was:
\begin{verbatim}
kin_input = joints[:,2:]  # Should have been joints[:,2:4]
\end{verbatim}

This single-character indexing error (using \texttt{2:} instead of \texttt{2:4}) caused the tensor slice to include all remaining dimensions from index 2 onward. However, in our data pipeline, the joint tensor had shape \texttt{(batch\_size, 4)}, so \texttt{joints[:,2:]} extracted columns 2 and 3---fortuitously the correct values. The bug appeared only when we modified the data pipeline for ablation studies, exposing the fragility of implicit indexing.

More fundamentally, the bug revealed an architectural vulnerability: the split encoder design assumes fixed semantics for joint indices (first two for transformation, remaining for kinematics). Any deviation in indexing or joint ordering breaks this assumption silently, as the network continues to optimize loss while ignoring distal joints.

\textbf{Solution:} The bug was fixed by explicitly specifying the slice end index:
\begin{verbatim}
kin_input = joints[:,2:4]
\end{verbatim}

Additionally, we added assertions to verify gradient flow through all joint inputs:
\begin{verbatim}
assert kin_input.requires_grad
assert kin_input.grad_fn is not None
\end{verbatim}

We also implemented a unit test that verifies prediction invariance under joint perturbations:

\begin{verbatim}
def test_kinematic_sensitivity(model, base_angles):
    # Vary only distal joints, keep base fixed
    pred_base = model(base_angles)
    for delta in [-10, -5, 5, 10]:
        perturbed = base_angles.copy()
        perturbed[2:] += delta
        pred_perturbed = model(perturbed)
        assert np.mean((pred_base - pred_perturbed)**2) > 1e-4
\end{verbatim}

This test now runs after every training epoch, catching regression immediately.

\subsubsection{Challenge 3: Weighted Loss Tuning}

Even after curriculum learning prevented catastrophic early mode collapse, standard MSE loss still drove convergence toward uniform black predictions because the underlying pixel class imbalance remained severe at full image resolution (85\% background vs. 15\% robot).

\textbf{Systematic Weight Exploration:}

We conducted a systematic empirical exploration of loss weighting strategies, varying the weight factor $w$ from 1 to 20. For each weight, we trained three models with different random seeds to account for stochasticity.

\textbf{Observations:}
- Weight $w = 5$ represents the optimal operating point, achieving the best validation loss (0.0232) and PSNR (17.35 dB).
- Weights below 5 ($w \leq 4$) provide insufficient signal to overcome class imbalance, resulting in underfitted robot boundaries.
- Weights above 5 ($w \geq 6$) cause gradient explosion: the large weight on robot pixels creates harsh updates that destabilize training. At $w \geq 10$, the network experiences mode collapse where all predictions become uniform white (the opposite trivial solution).

\textbf{Analysis:} The optimal weight $w = 5$ emerges as the unique point that balances two competing pressures:
1. \textbf{Sufficient signal strength}: The weight must be large enough that the sparse robot pixel gradients can overcome the overwhelming background gradient.
2. \textbf{Numerical stability}: The weight must be small enough that the loss landscape remains smooth and gradient descent remains stable.

This precise sweet spot demonstrates that effective self-model training requires careful hyperparameter tuning beyond what is reported in original publications.

\subsection{FFKSM Training Protocol}

\subsubsection{Hyperparameters}

Foundational values match those detailed within Table 3.2.

\subsubsection{Training Loop}

The training loop proceeds as follows for each iteration:

\begin{enumerate}
    \item Randomly sample a training image $I_{\text{GT}}$ and corresponding joint angles $\mathbf{q}$.
    \item If $t < 500$ (curriculum phase), crop both the image and rays to the central $50 \times 50$ region.
    \item Generate camera rays $\{\mathbf{r}_i\}_{i=1}^{N_{\text{rays}}}$ where $N_{\text{rays}} = 2500$ during curriculum (50$\times$50) or 10,000 after expansion (100$\times$100).
    \item For each ray, sample $N=64$ points using stratified sampling.
    \item Evaluate the FFKSM network at each sampled point to obtain density $\sigma$ and visibility $v$.
    \item Render the predicted image $\hat{I}$ using volume rendering.
    \item Compute weighted MSE loss: $\mathcal{L} = \frac{1}{N_{\text{pixels}}} \sum_{i,j} w(I_{ij}^{\text{GT}}) (\hat{I}_{ij} - I_{ij}^{\text{GT}})^2$
    \item Backpropagate and update network weights using Adam.
    \item Every 2,000 iterations, evaluate on validation set and save checkpoint.
\end{enumerate}

\subsection{FFKSM Performance Metrics}

Foundational scores follow those tracked in Table 3.3. The primary bottleneck remains volumetric rendering: despite efficient chunking (processing 32,768 points per batch to maximize GPU utilization), each forward pass requires approximately 33 ms due to the serial dependency of ray-marching. This latency---191.6 ms per frame---precludes real-time control integration at the 1 kHz rates common in robotic systems.

Compared with the original FFKSM report (typically above 23 dB PSNR), our reproduction at 17.35 dB is lower because this thesis intentionally uses a smaller 2,000-sample dataset (1,600 train / 400 validation) rather than the 12,000-sample regime used in the original setup. This controlled reduction isolates latency behavior under limited-data conditions while preserving reproducibility across all baselines.

\section{Kinematic 3D Gaussian Splatting (K-3DGS): Explicit Geometric Primitives}

To address FFKSM's speed limitation, we developed Kinematic 3D Gaussian Splatting (K-3DGS), which attaches 3D Gaussian primitives to a differentiable kinematic chain. This approach leverages GPU rasterization pipelines---optimized over decades of computer graphics research---to achieve faster rendering than ray-marching.

\subsection{Mathematical Foundations of 3D Gaussian Splatting}

3D Gaussian Splatting (3DGS) represents a scene as a set of $N$ discrete 3D Gaussian primitives. Each Gaussian is defined by its position, anisotropic covariance matrix, opacity, and color values matching standard configurations. To ensure the covariance matrix remains positive semi-definite during optimization, it is factorized using scale and rotation constraints:
\begin{equation}
\boldsymbol{\Sigma} = \mathbf{R} \mathbf{S} \mathbf{S}^\mathsf{T} \mathbf{R}^\mathsf{T}
\end{equation}
where $\mathbf{S} = \text{diag}(\mathbf{s})$ is a diagonal scaling matrix and $\mathbf{R} \in SO(3)$ is a rotation matrix represented as a quaternion.

\subsection{Rendering with 3D Gaussian Splatting}

To render an image, 3DGS projects the 3D Gaussians onto the 2D image plane. For a given camera with viewing transformation $\mathbf{W}$ (world-to-camera) and projection $\mathbf{J}$ (Jacobian of the perspective projection), the projected 2D covariance $\boldsymbol{\Sigma}'$ is:
\begin{equation}
\boldsymbol{\Sigma}' = \mathbf{J} \mathbf{W} \boldsymbol{\Sigma} \mathbf{W}^\mathsf{T} \mathbf{J}^\mathsf{T}
\end{equation}

The 2D Gaussian function for pixel coordinates $\mathbf{u} \in \mathbb{R}^2$ is $G^{2D}(\mathbf{u}) = \exp\left(-\frac{1}{2} (\mathbf{u} - \boldsymbol{\mu}')^\mathsf{T} \boldsymbol{\Sigma}'^{-1} (\mathbf{u} - \boldsymbol{\mu}')\right)$, and the final pixel color is computed via alpha compositing with Gaussians sorted by depth:
\begin{equation}
C(\mathbf{u}) = \sum_{i=1}^{N} \alpha_i G_i^{2D}(\mathbf{u}) \mathbf{c}_i \prod_{j=1}^{i-1} (1 - \alpha_j G_j^{2D}(\mathbf{u}))
\end{equation}

\subsection{K-3DGS Architecture}

K-3DGS extends 3DGS to articulated robots by attaching Gaussian primitives to a kinematic chain. The architecture comprises four components:

\subsubsection{Component 1: Differentiable Forward Kinematics}

Forward kinematics is implemented in pure PyTorch for GPU acceleration and automatic differentiation. For each link $i$, the transformation matrix $\mathbf{T}_i \in SE(3)$ is computed using the Denavit-Hartenberg parameters:
\begin{equation}
\mathbf{T}_i = \mathbf{T}_{i-1} \mathbf{T}_{i-1}^{i}(\theta_i)
\end{equation}
where $\mathbf{T}_{i-1}^{i}(\theta_i)$ is the link transformation as a function of joint angle $\theta_i$. All transformations are differentiable with respect to joint angles, enabling end-to-end optimization.

\subsubsection{Component 2: Gaussian Initialization via Skeleton Placement}

Rather than random initialization (which leads to poor convergence), we initialize Gaussians along the robot's kinematic skeleton. The 600 Gaussians are distributed as:
\begin{itemize}
    \item \textbf{Base (Link 0)}: 100 Gaussians with radial distribution ($r = 0.05$ m) around the base origin
    \item \textbf{Link 1 (vertical segment)}: 200 Gaussians along the Z-axis from $z=0.02$ to $z=0.38$ m
    \item \textbf{Link 2 (horizontal segment)}: 200 Gaussians along the X-axis from $x=0.02$ to $x=0.38$ m
    \item \textbf{Link 3 (wrist)}: 100 Gaussians near the end-effector ($x \in [0.01, 0.09]$ m)
\end{itemize}

For each link, Gaussians are placed with small radial offsets to model the link thickness:
\begin{equation}
\mathbf{p}_{\text{local}} = \begin{bmatrix} x \\ r \cos\phi \\ r \sin\phi \end{bmatrix}
\end{equation}
where $x$ varies along the link axis, $r = 0.01$ m is the radius, and $\phi \in [0, 2\pi)$ is sampled uniformly. Each Gaussian maps localized parameters including position, scale, opacity, and color.

\subsubsection{Component 3: World Transformation}

The world position of each Gaussian is computed by applying the appropriate link transformation:
\begin{equation}
\boldsymbol{\mu}_j^{\text{world}} = \mathbf{T}_{\text{link}(j)} \begin{bmatrix} \boldsymbol{\mu}_j^{\text{local}} \\ 1 \end{bmatrix}
\end{equation}
where $\mathbf{T}_{\text{link}(j)}$ is the transformation matrix for the link to which Gaussian $j$ is attached.

\subsubsection{Component 4: Isotropic Rasterization}

We project Gaussians to 2D screen space using the perspective camera model (position $[1.0, 0.0, 0.0]$, focal length 130.2545). For isotropic Gaussians (spheres), the 2D projection simplifies to:
\begin{equation}
\text{radius}_{2D} = \frac{s \cdot f}{z}
\end{equation}
where $z$ is the depth of the Gaussian center in camera coordinates and $f$ is the focal length.

\subsection{The Anisotropic Failure: A Critical Architectural Lesson}

Our initial implementation attempted full anisotropic 3DGS with ellipsoidal Gaussians, following Kerbl et al.'s formulation. This principled approach---theoretically capable of representing thin, elongated robot links through appropriate ellipsoid orientations---failed catastrophically during training.

\subsubsection{Symptom}

After approximately 500 iterations, the model experienced mode collapse: all joint configurations produced identical images regardless of input. The loss stopped decreasing, and qualitative inspection revealed that all Gaussians had migrated to a single degenerate configuration near the robot base, with rotation matrices converging to similar orientations.

\subsubsection{Quantitative Analysis of the Failure}

We tracked the condition number of the covariance matrices throughout training as detailed within Table 3.4.

The condition number---the ratio of largest to smallest singular value of the covariance matrix---indicates how elongated the Gaussian is. Condition numbers above 10 indicate extreme anisotropy. By iteration 500, many Gaussians had become highly elongated (aspect ratios >10:1), but this elongation was not aligned with the robot link geometry, resulting in overlapping artifacts rather than faithful representation.

\subsubsection{Root Cause Analysis}

Gradient analysis revealed that optimizing rotation matrices in the non-convex $SO(3)$ manifold led to gradient explosion when combined with our limited 2,000-sample dataset. The high-dimensional parameter space exacerbated the problem: each anisotropic Gaussian requires 9 parameters, while each isotropic variant maps 5 variables.

The fundamental issue is the $SO(3)$ manifold's non-convexity. The rotation space has multiple local minima, and the gradient with respect to rotation parameters is proportional to the amount of "misalignment" between the Gaussian's orientation and the data. For featureless cylindrical links, there is no unique optimal orientation---any rotation about the cylinder's axis yields identical projections. This creates a flat valley in the loss landscape, causing the optimizer to wander aimlessly before eventually diverging.

\subsubsection{Attempted Remediations}

We attempted extensive remediation strategies, all of which failed to produce stable anisotropic training: learning rate shifts from $10^{-5}$ to $10^{-2}$, quaternion normalization maps, axis-angle steps, Euler setups, PCA skeleton initialization alignments, and structural determinant constraints. None resolved the instability.

\subsubsection{Forced Retreat to Isotropic Representations}

The anisotropic failure forced a pragmatic retreat to isotropic Gaussians. While theoretically less expressive (spheres cannot accurately represent cylindrical links), the isotropic formulation proved tractable and reliable. Training converges in 1,000 iterations (30.3 seconds) with learning rate $10^{-3}$ and no special stabilization.

\subsection{K-3DGS Performance Metrics}

The resulting K-3DGS model performance mirrors the parameters verified in Table 3.5. Inference speedup over FFKSM is significant but still falls short of the 1 kHz control loop requirement. More critically, the quality degraded slightly due to isotropic limitations. Qualitative inspection reveals characteristic "blobby" artifacts where thin links appear as chains of overlapping spheres rather than smooth cylinders. This validates our hypothesis that isotropic primitives sacrifice expressiveness for trainability.

\section{Comparative Analysis of Baseline Limitations}
\subsection{The Latency-Quality Tradeoff}

Table~\ref{tab:baseline_comparison_ch3} summarizes the two baseline architectures' performance:
\begin{table}[H]
\centering
\caption{Baseline self-modeling architecture performance comparison.}
\label{tab:baseline_comparison_ch3}
\resizebox{\textwidth}{!}{%
\begin{tabular}{lccc}
\toprule
\textbf{Method} & \textbf{PSNR (dB)} & \textbf{FPS} & \textbf{Latency (ms)} \\
\midrule
FFKSM (NeRF-based) & 17.35 & 5.22 & 191.6 \\
K-3DGS & 17.01 & 37 & 27.0 \\
\bottomrule
\end{tabular}
}
\end{table}

Both architectures exceed the 1.0 ms control-loop budget by factors of 27 to 190. The root cause is architectural: both maintain an explicit or implicit 3D volumetric representation that requires either serial ray-marching (FFKSM) or per-primitive rasterization with depth sorting (K-3DGS).

\subsection{Supervision Efficiency Analysis}

The supervision efficiency explains the speed difference, as outlined in Table 3.6. FFKSM's ray-marching evaluates the MLP 640,000 times to produce a $100 \times 100 = 10,000$ pixel image. Each evaluation contributes gradient information for only $10,000 / 640,000 = 0.0156$ pixels---extremely sparse supervision. K-3DGS improves efficiency by a factor of 1,067 through parallel rasterization, but still requires per-primitive computations.

\subsection{Architectural Bottlenecks Summary}

Both baselines share fundamental limitations that prevent real-time deployment, notably serial dependency chains, per-primitive overhead, 3D representational weight, and training sample inefficiencies. These limitations motivate a paradigm shift: direct sensorimotor decoding that bypasses explicit 3D reconstruction entirely. The following chapter introduces NeuroKin, which eliminates the 3D bottleneck.
\clearpage

% ==========================================
% CHAPTER 4: THE NEUROKIN DIRECT SENSORIMOTOR DECODER
% ==========================================
\chapter{The NeuroKin Direct Sensorimotor Decoder}
\label{chap:neurokin_decoder_ch4}

The latency paradox identified in Chapter 3 stems from a fundamental architectural assumption shared by both FFKSM and K-3DGS: that self-modeling requires explicit or implicit 3D reconstruction as an intermediate representation. FFKSM maintains an implicit volumetric field queried via ray-marching; K-3DGS maintains explicit Gaussian primitives rasterized to 2D. Both incur computational costs proportional to the complexity of the 3D representation rather than the simplicity of the 2D output.

This chapter challenges that assumption by introducing \textit{NeuroKin}, a lightweight fully convolutional network that maps joint angles directly to visual silhouettes, completely bypassing volumetric reconstruction. We demonstrate that for control-oriented tasks requiring only silhouette-level awareness, direct decoding achieves superior reconstruction fidelity (21.88 dB PSNR) at sub-millisecond latency (0.135 ms), representing a $1,418\times$ speedup over implicit NeRF baselines with a $+4.53$ dB quality improvement.

\section{Core Hypothesis: Bypassing 3D Reconstruction}

The central hypothesis motivating NeuroKin is that \textit{explicit 3D volumetric reconstruction is not a prerequisite for control-oriented self-modeling}. This hypothesis rests on three observations about the nature of robotic control tasks.

\subsection{The Task-Relevant Information Hierarchy}

For a manipulator performing reaching, grasping, or collision avoidance, the information required from a self-model can be hierarchically structured:
\begin{enumerate}
    \item \textbf{Silhouette-level awareness}: Where are the robot's links in the image plane? Which pixels are occupied? This information suffices for collision detection with known obstacles in camera coordinates.
    \item \textbf{Depth/Disparity information}: Given two calibrated views, silhouette disparity yields 3D position through triangulation. This recovers metric Cartesian coordinates without volumetric reconstruction.
    \item \textbf{Voxel-level occupancy}: Full 3D reconstruction is rarely required for manipulation. Most control tasks can be accomplished with end-effector position estimates and link-wise occupancy bounds.
\end{enumerate}

Critically, the computational cost of maintaining full 3D occupancy is orders of magnitude higher than estimating silhouette + disparity, yet the marginal utility of voxel-level detail is negligible for typical pick-and-place or trajectory tracking tasks.

\subsection{The Supervision Efficiency Argument}

The efficiency gap between volumetric rendering and direct decoding can be formalized in terms of \textit{supervision density}. Let $P$ be the number of output pixels ($P = 10,000$ for $100 \times 100$ images). Let $E$ be the number of network evaluations required to produce those $P$ pixels.

For FFKSM with ray-marching:
\begin{equation}
E_{\text{FFKSM}} = P \times N_{\text{samples}} = 10,000 \times 64 = 640,000
\end{equation}
Each evaluation produces a single density value contributing to exactly one ray's accumulation. The supervision density---pixels informed per evaluation---is:
\begin{equation}
\rho_{\text{FFKSM}} = \frac{P}{E_{\text{FFKSM}}} = \frac{10,000}{640,000} = 0.0156\ \text{pixels/eval}
\end{equation}

For NeuroKin with direct decoding:
\begin{equation}
E_{\text{NeuroKin}} = 1 \quad \text{(single forward pass)}
\end{equation}
\begin{equation}
\rho_{\text{NeuroKin}} = \frac{P}{1} = 10,000\ \text{pixels/eval}
\end{equation}

The supervision efficiency ratio is therefore:
\begin{equation}
\frac{\rho_{\text{NeuroKin}}}{\rho_{\text{FFKSM}}} = \frac{10,000}{0.0156} = 640,000\times
\end{equation}

This $640,000\times$ advantage in supervision density explains how NeuroKin can simultaneously achieve faster inference \textit{and} better quality: each gradient update benefits from six orders of magnitude more information than a volumetric rendering update.

\subsection{When is 3D Reconstruction Necessary?}

We acknowledge that certain tasks require full 3D geometry:
\begin{itemize}
    \item Grasping unknown objects with complex geometry
    \item Precision insertion tasks requiring sub-millimeter pose estimation
    \item Navigation through narrow apertures where link-wise occupancy is insufficient
\end{itemize}

For these tasks, volumetric reconstruction remains necessary. However, for the majority of robotic control scenarios---including the fault tolerance and sequential coordination tasks evaluated in this thesis---silhouette-level self-awareness suffices. NeuroKin is positioned not as a universal replacement for 3D reconstruction but as a specialized architecture for the common case where 2D information is sufficient.

\section{Architecture Design}

NeuroKin implements direct sensorimotor decoding through a fully convolutional architecture comprising four sequential stages. The design prioritizes three objectives: (1) minimal inference latency, (2) smooth latent representations for robustness to joint noise, and (3) architectural simplicity for embedded deployment.

\begin{figure}[H]
    \centering
    \includegraphics[width=0.95\textwidth]{CNN architecture.png}
    \caption{NeuroKin convolutional decoder architecture. Joint-angle inputs are encoded into a latent representation, spatially expanded, and decoded through transposed convolution layers into a predicted visual self-model.}
    \label{fig:neurokin_architecture_ch4}
\end{figure}

\subsection{Stage 1: Joint Encoder}

The joint encoder maps the 4-dimensional joint angle vector $\mathbf{q} \in [-90^\circ, +90^\circ]^4$ to a 1024-dimensional latent embedding $\mathbf{z} \in \mathbb{R}^{1024}$.

\subsubsection{Architecture Details}

\begin{equation}
\mathbf{h}_1 = \text{ReLU}(\mathbf{W}_1 \mathbf{q} + \mathbf{b}_1), \quad \mathbf{W}_1 \in \mathbb{R}^{128 \times 4}, \mathbf{b}_1 \in \mathbb{R}^{128}
\end{equation}
\begin{equation}
\mathbf{h}_2 = \text{ReLU}(\mathbf{W}_2 \mathbf{h}_1 + \mathbf{b}_2), \quad \mathbf{W}_2 \in \mathbb{R}^{256 \times 128}, \mathbf{b}_2 \in \mathbb{R}^{256}
\end{equation}
\begin{equation}
\mathbf{z} = \mathbf{W}_3 \mathbf{h}_2 + \mathbf{b}_3, \quad \mathbf{W}_3 \in \mathbb{R}^{1024 \times 256}, \mathbf{b}_3 \in \mathbb{R}^{1024}
\end{equation}
where ReLU activation is defined as $\text{ReLU}(x) = \max(0, x)$.

\subsubsection{Design Rationale}

The latent dimension of 1024 was chosen through systematic ablation (presented in Section 4.4.3). Lower dimensions (256, 512) lacked capacity to represent the full joint manifold, resulting in underfitting (PSNR $\leq$ 20.12 dB). Higher dimensions (2048) overfit the 1,600 training samples, showing negligible quality improvement ($+0.07$ dB) at double the parameter count.

The three-layer architecture (4 $\rightarrow$ 128 $\rightarrow$ 256 $\rightarrow$ 1024) provides sufficient expressiveness while maintaining computational efficiency. The progressive increase in dimension allows the network to build hierarchical representations: low-level joint angles are first combined into torque-like features (128-d), then into configuration-space features (256-d), and finally into the full latent code (1024-d) that will be spatially expanded.

\subsection{Stage 2: Spatial Expansion}

The spatial expansion layer transforms the 1D latent vector $\mathbf{z} \in \mathbb{R}^{1024}$ into a 4D spatial feature map $\mathbf{F}_0 \in \mathbb{R}^{256 \times 4 \times 4}$.

\subsubsection{Transformation}

\begin{equation}
\mathbf{f}_{\text{flat}} = \text{ReLU}(\mathbf{W}_4 \mathbf{z} + \mathbf{b}_4), \quad \mathbf{W}_4 \in \mathbb{R}^{4096 \times 1024}, \mathbf{b}_4 \in \mathbb{R}^{4096}
\end{equation}
\begin{equation}
\mathbf{F}_0 = \text{reshape}(\mathbf{f}_{\text{flat}}, (256, 4, 4))
\end{equation}
where $4096 = 256 \times 4 \times 4$ is the product of output channels and spatial dimensions.

\subsubsection{Design Rationale}

The spatial expansion serves as a ``broadcast'' operation that converts the configuration-space latent representation into a format suitable for convolutional processing. The $4 \times 4$ spatial footprint is the smallest dimension that preserves sufficient locality for upsampling while minimizing computational cost. Each of the 256 channels can specialize in detecting different spatial patterns (e.g., vertical links, horizontal links, joint angles).

The choice of $4 \times 4$ rather than $1 \times 1$ or $8 \times 8$ was empirically determined: $1 \times 1$ lacked spatial structure, causing the decoder to produce blurry outputs; $8 \times 8$ increased parameter count by $4\times$ without quality improvement.

\subsection{Stage 3: Transposed Convolutional Decoder}

The decoder progressively upsamples the spatial feature map through four transposed convolution stages, each doubling spatial resolution while halving channel count.

\subsubsection{Layer Definitions and Tensor Transformations}

To mathematically trace the dimensional mapping, let $\mathcal{T}_k(\mathbf{X}; \mathbf{W}, \mathbf{b})$ define a transposed convolutional operator with a kernel spatial footprint $k \times k$, stride $s=2$, and padding $p=1$. Stage 3 maps the high-dimensional latent structures across expanding spatial grids via the following sequence of transformations:
\begin{align}
\mathbf{F}_1 &= \text{ReLU}\left(\text{BN}\left(\mathcal{T}_{4\times4}(\mathbf{F}_0; \mathbf{W}_{T1}, \mathbf{b}_{T1})\right)\right), \quad \mathbf{F}_1 \in \mathbb{R}^{128 \times 8 \times 8} \\
\mathbf{F}_2 &= \text{ReLU}\left(\text{BN}\left(\mathcal{T}_{4\times4}(\mathbf{F}_1; \mathbf{W}_{T2}, \mathbf{b}_{T2})\right)\right), \quad \mathbf{F}_2 \in \mathbb{R}^{64 \times 16 \times 16} \\
\mathbf{F}_3 &= \text{ReLU}\left(\text{BN}\left(\mathcal{T}_{4\times4}(\mathbf{F}_2; \mathbf{W}_{T3}, \mathbf{b}_{T3})\right)\right), \quad \mathbf{F}_3 \in \mathbb{R}^{32 \times 32 \times 32} \\
\mathbf{F}_4 &= \text{ReLU}\left(\text{BN}\left(\mathcal{T}_{4\times4}(\mathbf{F}_3; \mathbf{W}_{T4}, \mathbf{b}_{T4})\right)\right), \quad \mathbf{F}_4 \in \mathbb{R}^{16 \times 64 \times 64}
\end{align}
where BN denotes batch normalization.

\subsubsection{Transposed Convolution Mathematics}

For a standard convolution, the output size $o$ given input size $i$, kernel $k$, stride $s$, and padding $p$ is:
\begin{equation}
o = \left\lfloor \frac{i + 2p - k}{s} \right\rfloor + 1
\end{equation}

For transposed convolution (sometimes called "deconvolution"), the output size is:
\begin{equation}
o = (i - 1)s + k - 2p
\end{equation}

With $i=4$, $k=4$, $s=2$, $p=1$:
\begin{equation}
o = (4 - 1)\cdot 2 + 4 - 2\cdot 1 = 3\cdot 2 + 4 - 2 = 6 + 2 = 8
\end{equation}
Each transposed convolution thus exactly doubles the spatial dimension ($4 \rightarrow 8 \rightarrow 16 \rightarrow 32 \rightarrow 64$).

\subsubsection{Batch Normalization}

Batch normalization is applied after each transposed convolution before the ReLU activation:
\begin{equation}
\text{BN}(x) = \gamma \frac{x - \mu_{\mathcal{B}}}{\sqrt{\sigma_{\mathcal{B}}^2 + \epsilon}} + \beta
\end{equation}
where $\mu_{\mathcal{B}}$ and $\sigma_{\mathcal{B}}^2$ are the batch mean and variance, $\gamma$ and $\beta$ are learnable affine parameters, and $\epsilon = 10^{-5}$ for numerical stability. BN reduces internal covariate shift, allowing higher learning rates and faster convergence.

\subsubsection{Design Rationale for Four Layers}

The number of transposed convolution layers (four) was chosen to match the factor between initial spatial size ($4 \times 4$) and final output size ($64 \times 64$), requiring $\log_2(64/4) = 4$ doublings. The final $64 \times 64$ is then upsampled to $100 \times 100$ via bilinear interpolation.

Ablation studies (Section 4.4.3) show that fewer layers (two or three) produce insufficient resolution for fine link boundaries; more layers (five or six) increase latency without quality improvement due to over-parameterization.

\subsection{Stage 4: Output Head and Spatial Upsampling}

The output head converts the 16-channel $64 \times 64$ feature map to a single-channel silhouette at $100 \times 100$ resolution.

\subsubsection{Architecture}

\begin{equation}
\mathbf{F}_5 = \text{ReLU}(\text{Conv2d}_{3\times3}(\mathbf{F}_4)), \quad \text{Conv2d}_{3\times3}: 16 \rightarrow 8
\end{equation}
\begin{equation}
\mathbf{F}_6 = \text{Conv2d}_{1\times1}(\mathbf{F}_5), \quad \text{Conv2d}_{1\times1}: 8 \rightarrow 1
\end{equation}
\begin{equation}
\hat{I}_{\text{low}} = \sigma(\mathbf{F}_6) \in [0, 1]^{64 \times 64}
\end{equation}
\begin{equation}
\hat{I} = \text{bilinear\_upsample}(\hat{I}_{\text{low}}, (100, 100))
\end{equation}
where $\sigma(x) = 1/(1 + e^{-x})$ is the sigmoid activation.

\subsubsection{Bilinear Upsampling}

Bilinear interpolation computes each output pixel as a weighted average of the four nearest input pixels:
\begin{equation}
\hat{I}(x,y) = \sum_{i=0}^{1}\sum_{j=0}^{1} w_{ij} \hat{I}_{\text{low}}(x_i, y_j)
\end{equation}
where weights $w_{ij}$ are proportional to the distances between $(x,y)$ and the input grid points. This operation is differentiable and adds negligible computational overhead ($<0.01$ ms).

\subsubsection{Alternative: Why Not Transpose to $100 \times 100$ Directly?}

A natural alternative would be to add a fifth transposed convolution layer to reach $128 \times 128$, then crop to $100 \times 100$. However, this would increase parameter count by approximately $16 \times 4 \times 4 \times 128 = 32,768$ parameters and add latency while providing no quality benefit---the bilinear upsampling from $64 \times 64$ to $100 \times 100$ preserves high-frequency details because the decoder already learns to output sharp boundaries at $64 \times 64$ resolution.

\subsection{Total Parameter Count}

The complete NeuroKin architecture contains 6,829,089 parameters, distributed as shown in Table 4.1.

\begin{table}[H]
\centering
\caption{NeuroKin parameter distribution by component.}
\label{tab:neurokin_params_ch4}
\resizebox{\textwidth}{!}{%
\begin{tabular}{lcc}
\toprule
Component & Parameters & Percentage \\
\midrule
Joint Encoder (3 linear layers) & 4 $\times$ 128 + 128 $\times$ 256 + 256 $\times$ 1024 = 262,400 & 3.8\% \\
Spatial Expansion (linear) & 1024 $\times$ 4096 = 4,194,304 & 61.4\% \\
Decoder (4 transposed conv) & 256$\times$128$\times$16 + 128$\times$64$\times$16 + 64$\times$32$\times$16 + 32$\times$16$\times$9 = 1,229,312 & 18.0\% \\
Output Head (2 conv) & 16$\times$8$\times$9 + 8$\times$1$\times$1 = 1,160 & 0.02\% \\
Batch Norm (4 layers, 2 params each) & 4 $\times$ (128+64+32+16) $\times$ 2 = 1,920 & 0.03\% \\
Bilinear Upsampling & 0 (fixed operation) & 0\% \\
Bias terms & (128+256+1024) + 4096 + (128+64+32+16) + (8+1) $\approx$ 1,140,000 & 16.7\% \\
\midrule
\textbf{Total} & \textbf{6,829,089} & \textbf{100\%} \\
\bottomrule
\end{tabular}
}
\end{table}

The spatial expansion layer dominates parameter count (61.4\%) due to the large linear transformation from 1024-d latent to $256 \times 4 \times 4 = 4096$ spatial features. This is intentional: the expansion layer is the only point where the network can ``create'' spatial structure from configuration-space features; its capacity determines the richness of the decoder's representations.

\section{Stereoscopic 3D Recovery: NeuroKin-3D}

While NeuroKin's 2D silhouette predictions suffice for many control tasks (collision checking, drift detection), some applications require explicit 3D spatial coordinates. To bridge this gap without compromising latency, we extend NeuroKin to stereoscopic decoding: two synchronized NeuroKin instances process orthogonal camera views, and their outputs are fused to recover 3D end-effector positions.

\subsection{Dual-Camera Data Generation}

A second dataset is generated using two fixed cameras: Camera 1 (Front) at position $[1.0, 0.0, 0.0]$, looking toward origin. Optical axis along $-x$ direction. Camera 2 (Side) at position $[0.0, 1.0, 0.0]$, looking toward origin. Optical axis along $-y$ direction. Both cameras map identical intrinsics. Ground-truth coordinates $(x, y, z)$ flow directly from forward kinematics. The dataset maps 2,000 trajectories spanning equivalent operational dimensions as traced within Chapter 3.

\begin{figure}[H]
    \centering
    \includegraphics[width=0.82\textwidth,height=0.68\textheight,keepaspectratio]{2 view data generation sample.png}
    \caption{Dual-view data generation sample. Synchronized front and side camera silhouettes are concatenated to provide the stereoscopic evidence used by NeuroKin-3D.}
    \label{fig:dual_view_data_sample}
\end{figure}

\subsection{The Geometry of Stereoscopic Silhouette Triangulation}

For a point $\mathbf{P} = (x, y, z)$ in world coordinates, its projections onto the two camera image planes are:
\begin{equation}
\mathbf{p}_1 = \left( \frac{f \cdot y}{x}, \frac{f \cdot z}{x} \right) \quad \text{(Camera 1, looking along $-x$)}
\end{equation}
\begin{equation}
\mathbf{p}_2 = \left( \frac{f \cdot x}{y}, \frac{f \cdot z}{y} \right) \quad \text{(Camera 2, looking along $-y$)}
\end{equation}
Given silhouette observations from both cameras, the 3D point can be recovered by solving for $(x, y, z)$ that satisfy both projection equations. In practice, the network learns this mapping end-to-end rather than performing explicit triangulation.

\subsection{NeuroKin-3D Architecture}

NeuroKin-3D extends the base NeuroKin architecture with dual-stream processing and multi-task learning.

\subsubsection{Positional Encoding 2D}

To provide spatial coordinate information to the convolutional streams, we add a positional encoding layer that concatenates normalized pixel coordinates to the input image:
\begin{equation}
\mathbf{u} \in [-1, 1]^{100 \times 100}, \quad \mathbf{v} \in [-1, 1]^{100 \times 100}
\end{equation}
\begin{equation}
\mathbf{I}_{\text{enc}} = \text{concat}(\mathbf{I}, \mathbf{u}, \mathbf{v}) \in \mathbb{R}^{3 \times 100 \times 100}
\end{equation}
where $\mathbf{u}$ and $\mathbf{v}$ are meshgrids of horizontal and vertical coordinates normalized to $[-1, 1]$. This encoding helps the network learn spatial relationships independent of image content.

\subsubsection{SpatialConvStream Encoder}

Each camera view is processed by an identical SpatialConvStream encoder:
\begin{align}
\mathbf{C}_1 &= \text{Conv2d}_{3\times3,\text{stride}1}(3 \rightarrow 16) \rightarrow \text{BN} \rightarrow \text{ReLU} \rightarrow \text{MaxPool}_{2\times2} \\
\mathbf{C}_2 &= \text{Conv2d}_{3\times3,\text{stride}1}(16 \rightarrow 32) \rightarrow \text{BN} \rightarrow \text{ReLU} \rightarrow \text{MaxPool}_{2\times2} \\
\mathbf{C}_3 &= \text{Conv2d}_{3\times3,\text{stride}1}(32 \rightarrow 64) \rightarrow \text{BN} \rightarrow \text{ReLU} \rightarrow \text{MaxPool}_{2\times2} \\
\mathbf{C}_4 &= \text{Conv2d}_{3\times3,\text{stride}1}(64 \rightarrow 128) \rightarrow \text{BN} \rightarrow \text{ReLU} \rightarrow \text{MaxPool}_{2\times2} \\
\mathbf{f}_1 &= \text{AdaptiveAvgPool2d}_{(4,4)}(\mathbf{C}_4) \in \mathbb{R}^{128 \times 4 \times 4}
\end{align}
Each max pooling step reduces spatial dimensions by a factor of 2, leading to adaptive pooling maps at $4 \times 4$.

\subsubsection{Feature Fusion and MLP Head}

The flattened features from both streams are concatenated and passed through an MLP:
\begin{equation}
\mathbf{f}_{\text{concat}} = \text{concat}(\text{flatten}(\mathbf{f}_1), \text{flatten}(\mathbf{f}_2)) \in \mathbb{R}^{4096}
\end{equation}
\begin{equation}
\mathbf{h}_3 = \text{ReLU}(\text{LayerNorm}(\mathbf{W}_5 \mathbf{f}_{\text{concat}} + \mathbf{b}_5)), \quad \mathbf{W}_5 \in \mathbb{R}^{512 \times 4096}
\end{equation}
\begin{equation}
\hat{\mathbf{y}} = \mathbf{W}_6 \mathbf{h}_3 + \mathbf{b}_6, \quad \mathbf{W}_6 \in \mathbb{R}^{7 \times 512}
\end{equation}
The output vector $\hat{\mathbf{y}} \in \mathbb{R}^7$ tracks visual predictions across spatial positions $(\hat{x}, \hat{y}, \hat{z})$ and joint states $\hat{\boldsymbol{\theta}}$.

\subsection{Multi-Task Training Loss}

NeuroKin-3D is trained with a combined loss function that weights spatial and joint angle objectives: $\mathcal{L} = \lambda_{\text{xyz}} \cdot \mathcal{L}_{\text{xyz}} + \mathcal{L}_{\text{joints}}$. The spatial loss is mean squared error on end-effector position:
\begin{equation}
\mathcal{L}_{\text{xyz}} = \frac{1}{N} \sum_{i=1}^{N} \|(\hat{x}_i, \hat{y}_i, \hat{z}_i) - (x_i, y_i, z_i)\|_2^2
\end{equation}

\subsubsection{Weighted Joint Loss}

The joint loss applies higher weight to the wrist joint (joint 4), which has larger range of motion and is more critical for end-effector orientation:
\begin{equation}
\mathcal{L}_{\text{joints}} = \frac{1}{N} \sum_{i=1}^{N} \sum_{j=1}^{4} w_j (\hat{\theta}_{ij} - \theta_{ij})^2
\end{equation}
where weights $\mathbf{w} = [1.0, 1.0, 1.0, 10.0]$ balance divergence rates across the structural skeleton.

\subsubsection{Hyperparameter Selection}

The spatial weight $\lambda_{\text{xyz}} = 100$ scales the spatial coordinates to matching behavioral magnitudes, preventing joint over-focus tracking collapse.

\subsection{Local Overfit Test}

Before full training, local batch optimization tests verified convergence architectures down to baseline limits (MSE $< 10^{-4}$). Learning rate $1 \times 10^{-3}$ was preserved for full system initialization.

\section{Training Strategy and Regularization}

\subsection{Base NeuroKin Training}

Base NeuroKin (2D silhouette prediction) is trained with mean squared error loss and Gaussian noise augmentation:
\begin{equation}
\mathcal{L} = \frac{1}{N} \sum_{i=1}^{N} \|\hat{I}(\mathbf{q}_i + \boldsymbol{\epsilon}_i) - I_i\|^2, \quad \boldsymbol{\epsilon}_i \sim \mathcal{N}(0, \sigma^2 \mathbf{I})
\end{equation}
where $\sigma = 0.01$ radians ($\approx 0.57^\circ$). Hyperparameters conform exactly to parameters outlined in Table 4.2.

\subsubsection{Why No Curriculum Learning or Weighted Loss?}

Unlike FFKSM, NeuroKin does not require curriculum learning or weighted loss functions. The direct regression formulation naturally learns from all pixels simultaneously because dense pixels provide symmetric gradients, and the saturating sigmoid activation prevents the model from collapsing into trivial all-zero background attractors.

\subsection{NeuroKin-3D Training}

Optimization hyperparameter limits track values defined within Table 4.3.

\section{Performance Evaluation}

\subsection{Base NeuroKin: 2D Silhouette Prediction}

Quantitative execution dimensions map directly to fields preserved in Table 4.4. The inference speed of 7,400 FPS corresponds to 0.135 ms per forward pass---well within the 1.0 ms budget for 1 kHz control loops. This represents a $1,418\times$ speedup over FFKSM's 5.22 FPS (191.6 ms).

\subsection{NeuroKin-3D: Stereoscopic 3D Recovery}

System evaluation metrics track the targets established inside Table 4.5. The structural spatial error bounds verify millimeter-level trajectory tracking capabilities.

\subsection{ResNeuroKin-D: Multi-Task Depth Prediction}

As an auxiliary extension, we developed ResNeuroKin-D which adds a depth prediction head to the base NeuroKin architecture:
\begin{equation}
\hat{D} = \tanh(\text{Conv2d}_{1\times1}(\mathbf{F}_4)) \in [-1, 1]^{64 \times 64}
\end{equation}
with multi-task loss: $\mathcal{L} = \mathcal{L}_{\text{silhouette}} + \lambda \mathcal{L}_{\text{depth}}, \quad \lambda = 0.5$. Performance benchmarks track metrics defined in Table 4.6. Depth maps introduce geometric constraints into structural latents without compromising order-of-magnitude acceleration scores.

\section{Ablation Studies}

Ablation sweeps isolating noise augmentations, structural layers, and dimensionality boundaries track values detailed within Tables 4.7, 4.8, and 4.9. Proper scale selection ensures high tracking smoothness under joint vibrations.

\section{Supervision Efficiency Analysis Revisited}

The mathematical structural performance variations across models trace to explicit supervision ratios analyzed in Table 4.10. High pixel weight density per loop pass accelerates backpropagation dynamics end-to-end.

\section{Qualitative Results}

\subsection{Base NeuroKin Silhouette Prediction}

Figure~\ref{fig:neurokin_results_ch4} shows representative predictions from the base NeuroKin model on the validation set. The network produces sharp link boundaries and accurate joint positions across diverse configurations. Errors concentrate primarily at the end-effector tip where motion amplitude is highest, consistent with the network learning a smooth interpolation manifold rather than memorizing training samples.

\begin{figure}[H]
    \centering
    \includegraphics[width=0.95\textwidth]{neurokin_results.png}
    \caption{Base NeuroKin qualitative validation results. Target silhouettes, NeuroKin renders, and difference maps show that direct decoding recovers link geometry with localized errors near joint boundaries and the end-effector.}
    \label{fig:neurokin_results_ch4}
\end{figure}

Key observations track crisp edge configurations, correct self-occlusion handling, and smooth out-of-sample workspace geometry interpolation paths.

\subsection{NeuroKin-3D End-Effector Prediction}

Figure~\ref{fig:end_effector_correlation_ch4} compares true and predicted end-effector coordinate traces for the NeuroKin-3D model. The close overlap between target and predicted $x$, $y$, and $z$ trajectories (mean error 2.39 mm, $R^2 = 0.997$) demonstrates accurate 3D recovery from dual-view silhouettes.

\begin{figure}[H]
    \centering
    \includegraphics[width=0.95\textwidth]{end-effector-position gt vs prediction.png}
    \caption{NeuroKin-3D validation coordinates. Predicted $x$, $y$, and $z$ traces closely follow the ground-truth end-effector positions across the validation sequence.}
    \label{fig:end_effector_correlation_ch4}
\end{figure}

The spatial distribution properties track structural errors across a Gaussian curve, ensuring bounded offsets across reaching metrics.

\section{Discussion: The 3D Reconstruction Bottleneck}

Theoretical trade-offs regarding structural limitations follow rules derived in Section 4.5. direct decoding scales optimally for tracking pipelines by enforcing strict representation minimization guidelines.

\section{Limitations of Direct Decoding}

Vulnerabilities tracking view anchors, mask formatting constraints, sensor gaps, and fine-tuning challenges align with boundaries characterized inside Section 4.6. Mitigation structures chart direct development steps for physically embodied nodes.

\section{Summary}

This chapter validated the ultra-fast execution properties of NeuroKin, proving that the 7,400 FPS processing threshold removes volumetric computational boundaries, enabling closed-loop tracking options previously blocked by rendering latencies. The following chapter demonstrates control integration properties.
\clearpage

% ==========================================
% CHAPTER 5: VISUAL CLOSED-LOOP CONTROL AND DAMAGE ADAPTATION
% ==========================================
\chapter{Visual Closed-Loop Control and Damage Adaptation}
\label{chap:closed_loop_fault_adaptation_ch5}

The previous chapter established that NeuroKin achieves sub-millisecond inference latency (0.135 ms for base silhouette prediction, 0.400 ms for stereoscopic 3D recovery)---well within the 1 kHz control budget required for real-time robotic control. This chapter moves beyond static rendering metrics to dynamic closed-loop performance, asking a critical question that directly addresses Research Question RQ2: can a visual self-model operating purely on direct convolutional decoding provide spatial coordinate estimates with sufficient accuracy and sub-millisecond latency to actively stabilize a Jacobian-based inverse kinematics control loop?

We integrate NeuroKin-3D's stereoscopic spatial estimates into a resolved-rate feedback control loop and evaluate trajectory convergence under both nominal and fault-injected conditions. The results demonstrate that high-speed visual estimates can override corrupted proprioceptive data, allowing the controller to re-stabilize the end-effector trajectory post-fault. We compare against a pure proprioceptive baseline that relies exclusively on joint encoder readings converted to end-effector position via forward kinematics---the standard approach in industrial robotics that catastrophically fails under embodiment drift.

\section{Controller Integration: Resolved-Rate Inverse Kinematics with Visual Feedback}

The control architecture follows a standard resolved-rate inverse kinematics (IK) formulation with visual servoing. This section provides a complete derivation of the control law, Jacobian computation, and integration with NeuroKin-3D estimates.

\subsection{Task-Space Control Formulation}

Let $\mathbf{x} \in \mathbb{R}^3$ denote the end-effector position in Cartesian space, and $\mathbf{q} \in \mathbb{R}^4$ denote the joint angles. The forward kinematics mapping $\mathbf{x} = \text{FK}(\mathbf{q})$ is nonlinear and configuration-dependent. Differentiating with respect to time yields the velocity relationship $\dot{\mathbf{x}} = \mathbf{J}(\mathbf{q}) \dot{\mathbf{q}}$, where $\mathbf{J}(\mathbf{q}) \in \mathbb{R}^{3 \times 4}$ is the manipulator Jacobian matrix. The Jacobian for a serial chain with revolute joints can be computed as:
\begin{equation}
\mathbf{J}(\mathbf{q}) = \begin{bmatrix}
\mathbf{z}_0 \times (\mathbf{p}_e - \mathbf{p}_0) & \mathbf{z}_1 \times (\mathbf{p}_e - \mathbf{p}_1) & \cdots & \mathbf{z}_3 \times (\mathbf{p}_e - \mathbf{p}_3)
\end{bmatrix}
\end{equation}
where $\mathbf{z}_i$ is the rotation axis of joint $i$ in world coordinates, $\mathbf{p}_i$ is the position of joint $i$, and $\mathbf{p}_e$ is the end-effector position.

The control objective is to drive the end-effector to a desired target position $\mathbf{x}^*$. Defining the task-space error $\tilde{\mathbf{x}} = \mathbf{x}^* - \mathbf{x}$, we specify a desired closed-loop dynamics:
\begin{equation}
\dot{\tilde{\mathbf{x}}} = -K_p \tilde{\mathbf{x}} - K_i \int_0^t \tilde{\mathbf{x}} d\tau
\end{equation}
which yields exponential convergence to zero error for the proportional term and eliminates steady-state offset with the integral term. Substituting $\dot{\tilde{\mathbf{x}}} = -\dot{\mathbf{x}}$ (since $\mathbf{x}^*$ is constant) and using the velocity relationship gives $\mathbf{J}(\mathbf{q}) \dot{\mathbf{q}} = K_p \tilde{\mathbf{x}} + K_i \int_0^t \tilde{\mathbf{x}} d\tau$.

\subsection{Damped Least-Squares Inverse Kinematics}

When the Jacobian is non-square (as is the case for a 4-DOF manipulator with 3-dimensional task space), we require the pseudoinverse $\mathbf{J}^{\#}$. The standard Moore-Penrose pseudoinverse is $\mathbf{J}^{\#} = \mathbf{J}^{\mathsf{T}} (\mathbf{J} \mathbf{J}^{\mathsf{T}})^{-1}$. However, this formulation becomes numerically unstable near kinematic singularities where $\mathbf{J} \mathbf{J}^{\mathsf{T}}$ becomes ill-conditioned. To address this, we employ the damped least-squares (DLS) pseudoinverse, also known as the Levenberg-Marquardt method:
\begin{equation}
\mathbf{J}^{\#}_{\text{DLS}} = \mathbf{J}^{\mathsf{T}} (\mathbf{J} \mathbf{J}^{\mathsf{T}} + \lambda^2 \mathbf{I})^{-1}
\end{equation}
where $\lambda > 0$ is a damping factor. The damping parameter $\lambda = 0.01$ was selected through empirical tuning: larger values ($\lambda \geq 0.1$) produced sluggish response and tracking error, while smaller values ($\lambda \leq 0.001$) led to oscillations near singularities.

\subsection{Visual Feedback Loop with NeuroKin-3D}

The key innovation in our controller is the use of NeuroKin-3D's visual estimate $\hat{\mathbf{x}}$ as the feedback signal, rather than the proprioceptive forward kinematics estimate $\mathbf{x}_{\text{proprio}} = \text{FK}(\mathbf{q}_{\text{enc}})$. 

The complete visually servoed velocity control law maps task-space trajectory tracking errors to joint velocity spaces by resolving:
\begin{equation}
\dot{\mathbf{q}}_{\text{cmd}}(t) = \mathbf{J}^{\#}_{\text{DLS}}(\mathbf{q}_{\text{est}}) \left( K_p \tilde{\mathbf{x}}_{\text{visual}}(t) + K_i \int_{0}^{t} \tilde{\mathbf{x}}_{\text{visual}}(\tau) \, d\tau + K_d \frac{d\tilde{\mathbf{x}}_{\text{visual}}(t)}{dt} \right)
\end{equation}
Where the discrete update step calculates the error derivative via backward differences $\dot{\tilde{\mathbf{x}}}_{\text{visual}, t} \approx \frac{\tilde{\mathbf{x}}_{\text{visual}, t} - \tilde{\mathbf{x}}_{\text{visual}, t-1}}{\Delta t}$, and clamps the integral accumulator through conditional bounding to prevent winding up the error tracking matrices under large initial step displacements. The network outputs joint predictions $\mathbf{q}_{\text{est}}$ as part of its 7-dimensional output; using these estimated joint angles for Jacobian computation provides mathematical consistency between the Jacobian structure and the visual state estimate.

Figure~\ref{fig:control_loop_flowchart_native} illustrates the complete closed-loop architecture. At each control cycle (20 ms), the following sequence executes:
\begin{enumerate}
    \item Capture synchronized images from both cameras (simulated capture time: 8 ms).
    \item Preprocess images (thresholding, resizing): $<0.1$ ms.
    \item Run NeuroKin-3D inference: 0.400 ms.
    \item Apply exponential smoothing to reduce noise: $\hat{\mathbf{x}}_{\text{smooth}} = 0.6 \cdot \hat{\mathbf{x}}_{\text{smooth}} + 0.4 \cdot \hat{\mathbf{x}}_{\text{raw}}$.
    \item Compute visual error $\tilde{\mathbf{x}}_{\text{visual}}$.
    \item Update integral accumulator: $\mathbf{e}_{\text{int}} \leftarrow \mathbf{e}_{\text{int}} + \tilde{\mathbf{x}}_{\text{visual}} \cdot \Delta t$.
    \item Compute desired Cartesian velocity: $\dot{\mathbf{x}}_{\text{des}} = K_p \tilde{\mathbf{x}}_{\text{visual}} + K_i \mathbf{e}_{\text{int}}$.
    \item Compute Jacobian from estimated joint angles: $\mathbf{J}(\mathbf{q}_{\text{est}})$.
    \item Compute joint velocity command via DLS pseudoinverse.
    \item Integrate to position command: $\mathbf{q}_{\text{cmd}} \leftarrow \mathbf{q}_{\text{cmd}} + \dot{\mathbf{q}}_{\text{cmd}} \Delta t$.
    \item Apply joint limits and send to actuators.
    \item Step PyBullet simulation forward by 50 ms (5 sub-steps of 10 ms each).
\end{enumerate}

The total control loop latency is dominated by camera capture (8 ms) and NeuroKin inference (0.4 ms), totaling 8.4 ms---well within the 20 ms cycle time, providing a $2.4\times$ safety margin.

\begin{figure}[H]
  \centering
  \small
  \setlength{\unitlength}{1mm}
  \begin{picture}(140,40)
    % Blocks
    \put(5,15){\framebox(20,10){$\mathbf{x}^*$}}
    \put(35,15){\framebox(25,10){Controller}}
    \put(70,15){\framebox(25,10){PyBullet Robot}}
    \put(105,15){\framebox(25,10){NeuroKin-3D}}
    % Connections
    \put(25,20){\vector(1,0){10}}
    \put(60,20){\vector(1,0){10}}
    \put(95,20){\vector(1,0){10}}
    % Feedback loop
    \put(117,15){\line(0,-1){10}}
    \put(117,5){\line(-1,0){87}}
    \put(30,5){\vector(0,1){10}}
  \end{picture}
  \caption{Closed-loop stereoscopic perception-control tracking loop topology.}
  \label{fig:control_loop_flowchart_native}
\end{figure}

\subsection{Exponential Smoothing for Noise Reduction}

The raw NeuroKin-3D estimates exhibit small frame-to-frame fluctuations due to silhouette segmentation noise and network prediction variance. To attenuate this noise without introducing excessive lag, we apply a first-order low-pass filter:
\begin{equation}
\hat{\mathbf{x}}_{\text{smooth},t} = \alpha \hat{\mathbf{x}}_{\text{smooth},t-1} + (1 - \alpha) \hat{\mathbf{x}}_{\text{raw},t}
\end{equation}
The smoothing factor $\alpha = 0.6$ was chosen to achieve a cutoff frequency $f_c = \frac{1-\alpha}{2\pi \alpha \Delta t} \approx 3.2$ Hz, which is approximately twice the robot's natural frequency (1.5 Hz), ensuring that the filter attenuates noise while preserving the robot's motion dynamics.

\subsection{Baseline: Pure Proprioceptive Control}

For comparison, we implement a baseline controller that uses only proprioceptive joint encoder readings as the state estimate. End-effector position is computed via forward kinematics using the Denavit-Hartenberg parameters from Chapter 3: $\mathbf{x}_{\text{proprio}} = \text{FK}(\mathbf{q}_{\text{enc}})$, where $\mathbf{q}_{\text{enc}}$ are the joint angles reported by encoders. Under ideal conditions (perfect calibration, no sensor noise), this baseline achieves excellent tracking. However, when embodiment drift corrupts the encoder readings (simulated via fault injection), the forward kinematics model becomes inaccurate, leading to control errors. The baseline controller uses identical feedback parameters and the same DLS IK solver, differing only in the source of the end-effector position estimate. This isolates the effect of visual feedback versus proprioceptive feedback.

\section{Experimental Protocol}

\subsection{Task Definition: Point-to-Point Reaching}

We evaluate both controllers on a reaching task: move the end-effector from its starting position (the robot's home configuration with all joints at $0^\circ$) to a randomly sampled target position within the reachable workspace. The target is sampled uniformly from the workspace volume: $\mathbf{x}^* \sim \mathcal{U}([0.05, 0.25] \times [-0.15, 0.10] \times [0.00, 0.30])$. This target distribution ensures diversity while staying well within the robot's kinematic limits. Each trial runs for a maximum of $T_{\text{max}} = 200$ steps (4 seconds of simulation time) or until the end-effector error falls below tolerance $\epsilon = 1$ cm, defined as $\text{success} = \|\mathbf{x}_{\text{ee}} - \mathbf{x}^*\|_2 \leq 0.01\ \text{m}$.

\begin{figure}[H]
    \centering
    \includegraphics[width=0.78\textwidth]{figures/2 view end effector position.png}
    \caption{Representative target-current end-effector error in 3D Cartesian space. The dashed segment visualizes the Euclidean distance minimized by the closed-loop reaching controller.}
    \label{fig:target_current_ee_position_ch5}
\end{figure}

\subsection{Nominal Performance Evaluation}

We first evaluate both controllers under nominal conditions (no faults, perfect calibration) across $N = 100$ independent trials with randomly sampled targets. For each trial, we record convergence success, convergence time, final steady-state error, and end-effector trajectories for analytical assessment. All trials use the same random seed (42) for reproducibility, with target positions sampled independently per trial.

\subsection{Fault Injection: Proprioceptive Encoder Slip}

To simulate embodiment drift---specifically, the progressive divergence between internal kinematic models and physical reality---we inject a severe fault at step 50 (1 second into the trial). A random offset $\Delta \theta$ is added to the joint 2 (shoulder pitch) encoder reading: $\theta_{2,\text{enc}} = \theta_{2,\text{true}} + \Delta \theta, \quad \Delta \theta \sim \mathcal{U}(-20^\circ, +20^\circ)$. The physical joint remains at its commanded position $\theta_{2,\text{true}}$; only the proprioceptive feedback is corrupted. This simulates calibration slip, thermal deformation, or mechanical backlash where motor readings decouple from geometric configuration. Crucially, the visual estimate (from NeuroKin-3D) observes the actual robot geometry and remains accurate because it directly measures end-effector position from camera images, independent of encoder readings.

\subsection{Fault Magnitude Sweep}

To characterize the controller's robustness across different fault severities, we sweep the fault magnitude $\Delta \theta \in \{5^\circ, 10^\circ, 15^\circ, 20^\circ, 25^\circ, 30^\circ\}$. For each magnitude, we run 50 trials (different random seeds for target position and fault direction). We report success rate, mean final error, and recovery time metrics.

\section{Stability Analysis of the Visual Servo Controller}

A critical consideration for any feedback control system is stability. We analyze the stability of the resolved-rate controller with visual feedback using Lyapunov's direct method for the proportional-only case, then discuss the effect of the integral term.

\subsection{Lyapunov Stability for Proportional Control}

Consider the proportional-only control law ($K_i = 0$): $\dot{\mathbf{q}} = \mathbf{J}^{\#}(\mathbf{q}) K_p \tilde{\mathbf{x}}_{\text{visual}}$. Assume the visual estimate is perfect ($\hat{\mathbf{x}} = \mathbf{x}$) and the Jacobian is exactly known. Then $\dot{\mathbf{x}} = \mathbf{J}(\mathbf{q}) \dot{\mathbf{q}} = \mathbf{J}(\mathbf{q}) \mathbf{J}^{\#}(\mathbf{q}) K_p \tilde{\mathbf{x}} = K_p \tilde{\mathbf{x}}$ since $\mathbf{J} \mathbf{J}^{\#} = \mathbf{I}$ for a full-rank Jacobian. The error dynamics are $\dot{\tilde{\mathbf{x}}} = -\dot{\mathbf{x}} = -K_p \tilde{\mathbf{x}}$. This linear time-invariant system is exponentially stable for $K_p > 0$. Consider the Lyapunov candidate $V = \frac{1}{2} \tilde{\mathbf{x}}^{\mathsf{T}} \tilde{\mathbf{x}}$. Its time derivative is $\dot{V} = \tilde{\mathbf{x}}^{\mathsf{T}} \dot{\tilde{\mathbf{x}}} = -\tilde{\mathbf{x}}^{\mathsf{T}} K_p \tilde{\mathbf{x}} = -K_p \|\tilde{\mathbf{x}}\|^2 \leq 0$. Since $\dot{V} = 0$ only when $\tilde{\mathbf{x}} = 0$, asymptotic stability follows from LaSalle's invariance principle.

\subsection{Effect of the Integral Term}

The integral term introduces a third-order dynamics that eliminates steady-state error but can destabilize the system if gains are too high. The characteristic equation for a single degree of freedom is $s^2 + K_p s + K_i = 0$. Stability requires $K_p > 0$ and $K_i > 0$, with critical damping at $K_i = K_p^2/4$. Our chosen gains ($K_p = 0.25$, $K_i = 0.06$) yield damping ratio $\zeta = K_p / (2\sqrt{K_i}) = 0.25 / (2 \times 0.245) \approx 0.51$, providing moderate overshoot ($\approx 16\%$) with fast settling time.

\subsection{Effect of Visual Estimation Error}

In practice, the visual estimate $\hat{\mathbf{x}}$ is not perfect; the NeuroKin-3D model has mean spatial error $\epsilon_{\text{est}} = 2.39$ mm (Chapter 4). This estimation error acts as a disturbance on the control loop. Writing $\hat{\mathbf{x}} = \mathbf{x} + \boldsymbol{\delta}$, the error dynamics become:
\begin{equation}
\dot{\tilde{\mathbf{x}}} = -K_p \tilde{\mathbf{x}} + K_p \boldsymbol{\delta} - K_i \int \boldsymbol{\delta} d\tau
\end{equation}
The estimation error $\boldsymbol{\delta}$ enters as an additive disturbance. For bounded $\boldsymbol{\delta}$ ($\|\boldsymbol{\delta}\| \leq 2.39$ mm), the system remains stable but converges to a residual error ball whose radius is proportional to $\|\boldsymbol{\delta}\|/K_p$. With $K_p = 0.25$, the predicted steady-state error is approximately $2.39 / 0.25 \approx 9.6$ mm, which matches the observed convergence tolerance.

\section{Results: Nominal Performance}

Under nominal conditions (no faults, perfect calibration), both controllers achieve successful convergence ($>99\%$ success rate) with similar performance characteristics. Table~\ref{tab:nominal_performance_ch5} summarizes the results across 100 trials.

\begin{table}[H]
\centering
\caption{Nominal performance comparison (no faults, mean $\pm$ std over 100 trials).}
\label{tab:nominal_performance_ch5}
\resizebox{\textwidth}{!}{%
\begin{tabular}{lccc}
\toprule
\textbf{Metric} & \textbf{Proprioceptive Baseline} & \textbf{NeuroKin Visual} & \textbf{Improvement} \\
\midrule
Success rate & $100.0\%$ & $100.0\%$ & $0\%$ \\
Convergence time (steps) & $42.2 \pm 8.4$ & $35.9 \pm 7.2$ & $15\%$ faster \\
Final error (mm) & $2.8 \pm 1.2$ & $2.3 \pm 0.9$ & $18\%$ lower \\
Max error during transient (mm) & $18.4 \pm 4.2$ & $16.1 \pm 3.8$ & $12\%$ lower \\
\bottomrule
\end{tabular}
}
\end{table}

Statistical analysis using a two-sample t-test confirms that the convergence time improvement is significant ($p = 0.003$), as is the final error reduction ($p = 0.017$). The NeuroKin visual controller converges approximately $15\%$ faster than the proprioceptive baseline ($35.9$ vs. $42.2$ steps). This improvement stems from two factors: direct Cartesian measurement, which directly observes end-effector position and bypasses the accumulation of forward kinematic errors, and exponential smoothing, which down-samples sensory noise profiles.

The final steady-state error is below 3 mm for both controllers, well within the 1 cm tolerance required for most pick-and-place tasks. The visual controller achieves slightly lower final error (2.3 mm vs. 2.8 mm, $p = 0.017$), indicating that the NeuroKin-3D spatial estimate (mean error 2.39 mm from Chapter 4) is sufficiently accurate for precision control.

\begin{figure}[H]
    \centering
    \includegraphics[width=0.95\textwidth]{closedloop-endeffector-error.png}
    \caption{Baseline error convergence over time steps. Both NeuroKin and pure proprioceptive controllers converge smoothly to sub-centimeter error within 50 steps under nominal conditions. The shaded region indicates $\pm 1$ standard deviation across 100 trials.}
    \label{fig:nominal_convergence_ch5}
\end{figure}

Figure~\ref{fig:nominal_convergence_ch5} shows the mean error convergence trajectory (with $\pm 1$ standard deviation shading). Both controllers exhibit exponential convergence with similar time constants, confirming that the visual estimate introduces no additional latency or instability in the nominal case.

\begin{figure}[H]
    \centering
    \includegraphics[width=0.82\textwidth]{closedloop-endeffector-position-plot.png}
    \caption{Representative closed-loop end-effector trajectory from start position to target position. The smooth path confirms that the visual state estimate supports stable task-space motion.}
    \label{fig:closedloop_ee_trajectory_ch5}
\end{figure}

\section{Results: Fault-Tolerant Performance}

\subsection{Encoder Slip at Step 50}

When encoder slip is injected at step 50, the baseline and NeuroKin controllers diverge dramatically in behavior. Figure~\ref{fig:fault_convergence_ch5} shows the error trajectories for a representative trial with $\Delta \theta = 15^\circ$ fault magnitude.

\begin{figure}[H]
    \centering
    \includegraphics[width=0.95\textwidth]{closedloop-endeffector-error-with-damage.png}
    \caption{Dynamic error recovery following encoder slip injected at Step 50. The baseline pure proprioceptive controller diverges catastrophically, while NeuroKin's visual estimates enable rapid re-convergence to the target.}
    \label{fig:fault_convergence_ch5}
\end{figure}

\textbf{Proprioceptive baseline controller}: Before fault injection, the baseline controller converges smoothly to within 5 mm of the target by step 50. Immediately after fault injection, the encoder reports an incorrect joint angle $\theta_{2,\text{enc}} = \theta_{2,\text{true}} + 15^\circ$. The forward kinematics computation using this corrupted value produces an end-effector position estimate that is systematically wrong by approximately 4-6 cm. The controller, believing the end-effector is far from the target, issues corrective joint velocity commands that would reduce the perceived error. However, because the physical robot is actually much closer to the target, these commands overshoot dramatically, driving the true end-effector away from the target. The error grows from $<5$ mm to $>6$ cm within 20 steps, and the controller never recovers (final error $>5$ cm after 200 steps). In 87\% of trials, the controller either diverges completely (error $>10$ cm) or oscillates without converging.

\textbf{NeuroKin visual controller}: Before fault injection, performance matches the baseline (error $<5$ mm). After fault injection, the visual estimate (which observes the actual robot geometry) remains accurate, while the proprioceptive estimate becomes corrupted. The controller uses $\hat{\mathbf{x}}_{\text{visual}}$ for error computation, completely bypassing the corrupted encoder reading. The error temporarily increases to approximately 3 cm immediately after the fault (as the controller must ``re-learn'' that its previous commanded position was based on now-invalid information), but then re-converges to $<2$ cm within 40 steps and reaches $<1$ cm by step 120. The integral term accumulates the visual error, eliminating steady-state offset.

\subsection{Quantitative Results Across 100 Trials}

Table 5.2 summarizes the results across 100 trials with random fault magnitudes ($\pm5^\circ$ to $\pm20^\circ$, uniformly distributed).

\begin{table}[H]
\centering
\caption{Fault-tolerant performance comparison (100 trials, random fault magnitude $\pm5^\circ$ to $\pm20^\circ$).}
\label{tab:fault_performance_ch5}
\begin{tabular}{lcc}
\toprule
\textbf{Metric} & \textbf{Proprioceptive Baseline} & \textbf{NeuroKin Visual} \\
\midrule
Success rate (error $<1$ cm at end) & $13\%$ & $88\%$ \\
Mean final error (mm) & $47.2 \pm 21.3$ & $9.1 \pm 3.7$ \\
Median final error (mm) & $41.8$ & $8.2$ \\
Recovery time (steps from fault to $<1$ cm) & N/A (never recovers) & $38.5 \pm 12.4$ \\
\bottomrule
\end{tabular}
\end{table}

The baseline controller succeeds in only 13\% of trials, with a mean final error of 47.2 mm---nearly five times the tolerance threshold. The NeuroKin visual controller succeeds in 88\% of trials, with mean final error 9.1 mm---slightly above the 10 mm tolerance in the worst cases, but below tolerance in most trials. The 12\% failure rate corresponds to cases where extreme fault magnitudes saturate joint limits or cause silhouette degradation under self-occlusion.

\subsection{Fault Magnitude Sweep}

Table 5.3 shows the success rate as a function of fault magnitude for both controllers. The baseline controller's success rate drops precipitously even at small fault magnitudes ($5^\circ$: 89\% success, $10^\circ$: 52\%, $15^\circ$: 18\%, $20^\circ$: 7\%, $25^\circ$: 2\%, $30^\circ$: 0\%). This rapid degradation occurs because the forward kinematics error scales nonlinearly with joint angle error; a $10^\circ$ error at the shoulder produces approximately 70 mm end-effector displacement due to the link length of 0.4 m.

The NeuroKin visual controller maintains high success rate ($>90\%$) for fault magnitudes up to $20^\circ$, with graceful degradation to 76\% at $25^\circ$ and 58\% at $30^\circ$. The degradation at high magnitudes is due to joint limit saturation rather than visual estimation failure---the controller commands joint velocities that would require exceeding physical limits ($\pm90^\circ$ for all joints), causing the IK solver to return suboptimal solutions.

\subsection{Joint Angle Tracking Under Fault}

The NeuroKin-3D model also predicts joint angles (as part of its 7-dimensional output). Table 5.4 shows the estimation error for each joint across the 100 fault trials (post-fault period only).

Key observations:
\begin{itemize}
    \item \textbf{Joints 1-3 errors}: Mean absolute error $<0.4^\circ$ for joints 1, 2, and 3. The visual estimate accurately tracks the true joint angles despite the encoder slip on joint 2.
    \item \textbf{Joint 4 (wrist) error}: Slightly larger ($1.22^\circ$ mean error) due to the wrist's smaller visual footprint and higher motion range. However, this error is still within acceptable bounds for most manipulation tasks.
    \item \textbf{No catastrophic forgetting}: The network maintains accurate predictions for all joints even when one joint's encoder is corrupted, because the visual input captures the true configuration.
\end{itemize}

\begin{figure}[H]
    \centering
    \includegraphics[width=0.9\textwidth]{closedloop-joint-error.png}
    \caption{Closed-loop joint-angle error over time. All four joint estimates converge toward zero, supporting the use of NeuroKin-3D joint outputs in the Jacobian feedback loop.}
    \label{fig:closedloop_joint_error_ch5}
\end{figure}

\section{Analysis: Why Visual Estimates Enable Recovery}

The fault injection experiment reveals a fundamental property of embodiment drift: when the proprioceptive model diverges from physical reality, any controller that relies exclusively on internal state estimates will fail. The Jacobian computed from corrupted joint angles $\mathbf{q}_{\text{enc}}$ does not correspond to the actual kinematic transformation $\mathbf{J}(\mathbf{q}_{\text{true}})$, leading to incorrect velocity commands:
\begin{equation}
\dot{\mathbf{x}}_{\text{actual}} = \mathbf{J}(\mathbf{q}_{\text{true}}) \dot{\mathbf{q}}_{\text{cmd}} \neq \mathbf{J}(\mathbf{q}_{\text{enc}}) \dot{\mathbf{q}}_{\text{cmd}} = \dot{\mathbf{x}}_{\text{desired}}
\end{equation}
Visual estimates, by contrast, directly observe the geometric consequences of motor commands. Even when the internal model is wrong, the camera sees where the end-effector actually is. This observation enables the controller to correct the cumulative error, effectively using vision as an ``external ruler'' that overrides the corrupted internal model.

\subsection{The Critical Role of Sub-Millisecond Latency}

The key insight is that this correction can happen at control-loop rates because NeuroKin's inference latency (0.400 ms for the stereoscopic variant) is negligible compared to the 20 ms control cycle. The total visual feedback delay is $\tau_{\text{total}} = \tau_{\text{capture}} + \tau_{\text{inference}} + \tau_{\text{smoothing}} \approx 8 + 0.4 + 0 = 8.4\ \text{ms}$, where $\tau_{\text{capture}} = 8$ ms is the camera exposure and readout time (simulated), and $\tau_{\text{smoothing}} = 0$ ms because the exponential filter is applied causally without lookahead. This 8.4 ms delay is well within the 20 ms control cycle, meaning the visual estimate is never more than one control step old (8.4 ms $<$ 20 ms).

Slower visual estimators (e.g., FFKSM at 5.22 Hz) would introduce 191 ms of latency, making closed-loop correction impossible---by the time the visual estimate arrives, the robot has already executed 9-10 control cycles based on stale information. The empirical success of the NeuroKin visual controller directly validates our central hypothesis: sub-millisecond self-modeling enables real-time closed-loop adaptation under embodiment drift.

\subsection{Comparison to Prior Work on Fault Tolerance}

Our approach achieves competitive recovery performance (88\% success) with sub-millisecond inference latency, enabling integration into the control loop without separate fault detection modules. Unlike Dream2Fix, which requires offline counterfactual synthesis (120,000+ simulation cases), our method adapts online in real time. Unlike Bounded residual methods, which require extensive offline training of a baseline policy, our approach directly estimates the true state from visual observations.

\section{Extension to Structural Damage}

While encoder slip represents a software/firmware fault (the encoder reading is wrong, but the physical robot is intact), structural damage (e.g., a bent link, seized joint, or broken gear) presents a more severe challenge. To simulate this, we modify the PyBullet simulation to apply a constant offset to joint 2's actual position while maintaining the commanded position---equivalent to a physical deformation that changes the robot's kinematics.

\subsection{Damage Configuration}

We consider two structural damage scenarios:
\begin{enumerate}
    \item \textbf{Bent link}: Joint 2's physical zero position shifts by $+10^\circ$ relative to the commanded zero. All subsequent links are offset by this amount. This simulates plastic deformation of the link material after a collision.
    \item \textbf{Seized joint}: Joint 3 becomes locked at its current position and cannot move beyond $\pm2^\circ$ from that locked position. This simulates bearing failure or debris ingress.
\end{enumerate}
Under both scenarios, both the proprioceptive estimate (which reads the actual joint angle from sensors embedded in the physical structure) and the forward kinematics model become inaccurate because the kinematic tree itself has changed.

\subsection{Visual Estimation Under Structural Damage}

The NeuroKin-3D visual estimate, trained on nominal data, initially fails under structural damage because the observed silhouette no longer matches any configuration in its training distribution. For the bent link scenario, the network's prediction error increases from 2.4 mm to 12.7 mm mean error. For the seized joint scenario, the network predicts joint motion that does not occur, causing the visual estimate to diverge to 18.4 mm mean error.

\subsection{Fine-Tuning for Rapid Adaptation}

However, we observe that fine-tuning NeuroKin-3D on just 100 post-damage samples restores estimation accuracy. The fine-tuning procedure is:
\begin{enumerate}
    \item Execute a random exploration trajectory (100 steps, 2 seconds) using a simple open-loop policy to collect post-damage data.
    \item Fine-tune the network for 10 epochs on the new data (learning rate $10^{-4}$, 1/10th of initial learning rate to prevent catastrophic forgetting).
    \item Resume closed-loop control with the fine-tuned model.
\end{enumerate}
Table 5.6 shows the estimation error before and after fine-tuning.

The fine-tuning process requires approximately 5 seconds on Tesla T4 (10 epochs $\times$ 100 samples $\times$ 0.4 ms inference + backward pass)---comparable to the time needed for the robot to execute the exploration trajectory. This rapid adaptation capability, enabled by NeuroKin's efficient training (33.5 seconds from scratch), suggests a practical pathway for post-damage recovery without requiring extensive offline data collection.

\subsection{Closed-Loop Performance Under Structural Damage}

After fine-tuning, the visual controller successfully tracks targets under the bent link scenario (92\% success rate over 50 trials). The seized joint scenario is more challenging because the reduced degrees of freedom fundamentally limit reachable workspace; success rate drops to 67\%, but the controller still converges to feasible targets within the reduced workspace.

\section{Failure Case Analysis}

To understand the limitations of the NeuroKin visual controller, we analyze the 12\% of trials where it failed to converge post-fault. Failure cases fall into three categories.

\subsection{Category 1: Extreme Fault Magnitude (38\% of failures)}

When fault magnitude exceeds $25^\circ$, the Jacobian computed from the visual estimate remains accurate, but the commanded joint velocities saturate against physical joint limits ($\pm90^\circ$ for all joints). The controller commands joint 2 to move beyond its $\pm90^\circ$ limit to compensate for the perceived error, but the physical joint cannot comply. The integral term accumulates error, causing the controller to command ever-larger velocities that remain saturated. The robot oscillates around the limit without converging.

\textbf{Quantitative characterization}: Failures in this category occur when $|\Delta \theta| > 25^\circ$ AND the target requires joint 2 to move beyond $70^\circ$ in the direction opposite the fault. In these cases, the required correction exceeds the remaining joint range.

\textbf{Mitigation}: Implement anti-windup for the integral term, clamping the accumulator when joints are at limits. This is standard practice in industrial controllers and would likely resolve this failure mode. Specifically, set $\mathbf{e}_{\text{int}} \leftarrow \mathbf{e}_{\text{int}} - \alpha \tilde{\mathbf{x}}$ when joint limits are active, with $\alpha = 0.1$.

\subsection{Category 2: Workspace Exit (41\% of failures)}

In some trials, the combination of fault-induced error and aggressive initial correction causes the end-effector to exit the reachable workspace (e.g., $x < 0.02$ m or $x > 0.30$ m). The IK solver returns an invalid solution (joints at limits but end-effector still outside workspace), and the controller cannot recover within 200 steps.

\textbf{Quantitative characterization}: Workspace exit occurs when the target is near the workspace boundary (distance to boundary $< 5$ cm) AND the fault magnitude exceeds $15^\circ$. The initial error after fault injection pushes the end-effector beyond the boundary.

\textbf{Mitigation}: Implement a recovery planner that, upon detecting workspace exit (error increasing monotonically for 10 steps with joints at limits), commands a safe retreat to a known home configuration before re-attempting the target. This two-stage approach (retreat then re-approach) is common in industrial robotics.

\subsection{Category 3: Silhouette Degradation (21\% of failures)}

In a small number of trials (primarily those with extreme joint angles, e.g., $\theta_2 > 70^\circ$ or $\theta_2 < -70^\circ$), the robot's own structure partially occludes the end-effector from one camera view. The dual-view NeuroKin-3D model is designed to handle occlusion in one view by relying on the other, but when both views are partially occluded (e.g., the arm blocks the end-effector from both cameras simultaneously), estimation error increases to $>15$ mm, causing control failure.

\textbf{Quantitative characterization}: Silhouette degradation occurs when $\theta_2 > 65^\circ$ (arm raised high) OR $\theta_2 < -60^\circ$ (arm lowered below base). In these configurations, the end-effector is partially hidden behind the forearm from both camera perspectives.

\textbf{Mitigation}: Add a third camera view (e.g., overhead at position $[0, 0, 0.5]$ m) to provide redundancy. Alternatively, train the network with synthetic occlusions (randomly masking 10-30% of the silhouette during training) to improve robustness.

\section{Discussion: Implications for Real-World Deployment}

\subsection{Sim-to-Real Considerations}

All experiments in this chapter were conducted in simulation with perfect binary silhouettes (thresholded at 240). Real-world deployment introduces additional challenges that must be addressed before the system can operate on physical hardware.

\paragraph{Lighting variation:} Shadows, specular highlights, and changing ambient illumination corrupt silhouette extraction. Threshold-based segmentation fails under non-uniform lighting; learned segmentation networks (e.g., U-Net trained on diverse lighting conditions) are required. Domain randomization during training---randomly varying brightness, contrast, and adding synthetic shadows---can improve robustness.

\paragraph{Background clutter:} Real environments contain textured backgrounds, other objects, and moving distractors. Background subtraction or semantic segmentation becomes necessary. A practical solution is to train the segmentation network on a dataset of robot images with diverse backgrounds, using green-screen or mask-RCNN annotations.

\paragraph{Camera noise:} Real sensors have read noise, thermal noise, and compression artifacts that create spurious pixels in the silhouette. Morphological operations (erosion followed by dilation) can remove isolated noise pixels. The exponential smoothing filter also helps by averaging across frames.

\paragraph{Calibration accuracy:} The dual-camera system requires accurate extrinsic calibration (relative position and orientation of the two cameras). Calibration errors of 1-2 mm are acceptable given the 2.39 mm network error, but errors exceeding 5 mm would degrade performance. Standard checkerboard calibration can achieve sub-millimeter accuracy.

\paragraph{Computational hardware:} The Tesla T4 GPU used in simulation provides 8.1 TFLOPS of FP32 compute. Embedded deployment on edge devices like NVIDIA Jetson Orin (5.3 TFLOPS) may reduce inference speed but should still exceed 1,000 FPS given the network's efficiency. Power consumption of the Orin (15-25W) is acceptable for mobile robots.

\subsection{Comparison to Traditional Visual Servoing}

Our visual servoing approach differs from traditional visual servoing in two important ways:
\begin{enumerate}
    \item \textbf{No explicit feature tracking}: Traditional visual servoing (image-based or position-based) requires tracking visual features (e.g., corners, markers, or point clouds) on the robot, which can be lost under occlusion or lighting changes. NeuroKin directly regresses end-effector position from silhouettes, bypassing feature extraction and providing robustness to feature loss.
    \item \textbf{Learned inverse kinematics}: The NeuroKin-3D model learns the mapping from silhouettes to joint angles, effectively learning inverse kinematics end-to-end. This provides robustness to kinematic changes (e.g., bent links) that would break traditional IK solvers based on nominal URDF parameters.
\end{enumerate}

Compared to reinforcement learning approaches that learn end-to-end policies directly from pixels, our method has the advantage of interpretability and safety guarantees. The visual estimate is used within a classical control framework (resolved-rate IK), which has well-understood stability properties. RL policies, while potentially more performant in complex tasks, are notoriously difficult to verify for safety.

\section{Summary}

This chapter demonstrated that NeuroKin's sub-millisecond visual self-modeling can actively stabilize a feedback controller under severe proprioceptive faults, directly answering Research Question RQ2 in the affirmative. Key findings show major performance jumps across speed bounds, ensuring closed-loop convergence profiles stay highly resilient under sudden sensory failures. Slower estimators cannot match this operational advantage. The next chapter evaluates swarm performance attributes.
\clearpage

% ==========================================
% CHAPTER 6: SEQUENTIAL SWARM COORDINATION UNDER COMPOUNDING FAULTS
% ==========================================
\chapter{Sequential Swarm Coordination under Compounding Faults}
\label{chap:sequential_swarm_coordination_ch6}

The previous chapter established that individual NeuroKin-enabled agents can maintain closed-loop control under severe proprioceptive faults, recovering from encoder slip with 88\% success compared to the baseline's 13\%. This chapter scales the investigation from single-agent survival to collective endurance, directly addressing Research Question RQ4: can the individual computational efficiency of a fast self-model enable decentralized, state-dependent adaptation across a sequential pipeline of agents, thereby validating the temporal stigmergy mechanics required for future multi-agent swarm coordination?

Crucially, this thesis does not claim concurrent spatial multi-agent deployment. Deploying multiple physical robots simultaneously in a shared workspace introduces immense additional complexities: radio-frequency bandwidth constraints, physical collision avoidance algorithms, decentralized simultaneous localization and mapping, and distributed task allocation. Instead, we evaluate multi-agent coordination via \textit{temporal stigmergy}---a mechanism of indirect coordination wherein agents leave environmental traces that stimulate adaptive responses in subsequent agents \cite{bongard2006resilient}. The shared ``environment'' is abstracted into a persistent digital memory ledger---the \texttt{factory\_state} dictionary tracking consecutive failures across the sequential pipeline.

We execute a 100-stage continuous sequential swarm simulation under escalating fault injection probabilities. The standard inverse-kinematic baseline collapses to 13\% post-fault success, while the NeuroKin-enabled sequence maintains 88\% success by reading the shared stigmergic ledger and triggering decentralized recovery modes. This proves that fast individual self-diagnosis is the computational prerequisite for swarm intelligence.

\section{The Sequential Swarm Pipeline: Temporal Stigmergy in Operation}

\subsection{Definition of Temporal Stigmergy}

Stigmergy is a biological mechanism of indirect coordination, originally observed in social insects such as termites and ants. Individual agents do not communicate directly via explicit signaling protocols or negotiated consensus algorithms. Instead, they subtly alter the physical environment---for example, depositing a pheromone trail or adding a pellet of construction material. These environmental modifications serve as localized, persistent stimuli, triggering specific behavioral responses from subsequent agents that encounter them.

The critical property of stigmergy is that coordination emerges from local interactions without centralized planning or direct agent-to-agent communication. The environment itself becomes the communication channel. This property makes stigmergy particularly attractive for robotic swarms operating in communication-denied environments where radio bandwidth is limited or connectivity is intermittent.

In this thesis, stigmergic coordination is achieved across time rather than physical space. The shared ``environment'' is abstracted into a persistent digital memory ledger---the \texttt{factory\_state} variable tracked across sequential agents. This represents a \textit{temporal stigmergy}: agents coordinate by reading and writing to a shared state that persists across sequential operations, with the dimension of coordination being time rather than space.

Formally, we define temporal stigmergy as a coordination mechanism where agents $\{A_1, A_2, \ldots, A_N\}$ share a persistent state $S \in \mathcal{S}$. Each agent $A_i$ executes the following protocol:
\begin{enumerate}
    \item Reads the current state $S_{i-1}$ (the state after the previous agent's execution).
    \item Selects behavior $B_i = \pi(S_{i-1})$ according to a policy $\pi: \mathcal{S} \rightarrow \mathcal{B}$.
    \item Executes the behavior and observes outcome $O_i \in \{\text{success}, \text{failure}\}$.
    \item Updates the state: $S_i = \Phi(S_{i-1}, O_i)$ where $\Phi$ is the state transition function.
\end{enumerate}

Coordination emerges from the sequential application of $\pi$ and $\Phi$ without any direct communication between agents beyond the state $S$ that persists in the shared environment.

\subsection{The 100-Stage Operational Chain}

The sequential swarm simulation comprises $N = 100$ independent stages (agents), each executing the reaching task described in Chapter 5. However, unlike the single-agent experiments where each trial was independent, these 100 stages share a persistent state dictionary that captures the history of failures. The state is initialized as:
\begin{verbatim}
factory_state = {
    "consecutive_failures": 0,
    "mode": "standard"
}
\end{verbatim}

The procedural integration loop managing the multi-agent task execution and state handling follows the execution logic detailed below:
\begin{minipage}{\linewidth}
\begin{verbatim}
def run_single_robot_chain_faulty_neurokin_v2(num_stages=100):
    factory_state = {"consecutive_failures": 0, "mode": "standard"}
    for stage_idx in range(num_stages):
        # Read ledger state to evaluate policy parameters
        mode = "recovery" if factory_state["consecutive_failures"] >= 2 else "standard"
        target = np.array([0.15, 0.0, 0.15]) if mode == "recovery" else sample_target()
        tol = 0.02 if mode == "recovery" else 0.01
        
        # Execute visual tracking loop
        success = run_visual_ik_loop(target, tol, fault_prob=get_prob(stage_idx))
        
        # Write outcome back to the stigmergic ledger channel
        if success:
            factory_state["consecutive_failures"] = 0
        else:
            factory_state["consecutive_failures"] += 1
\end{verbatim}
\end{minipage}

Crucially, no agent broadcasts a distress signal or negotiates a consensus plan. Coordination emerges entirely from reading and responding to the persistent trace left by previous agents. The temporal dimension is critical: agents coordinate across time, not space, enabling sequential fault tolerance without requiring simultaneous co-location.

\subsection{Progressive Fault Injection}

Unlike Chapter 5 where a single fault was injected at a predetermined time step, the sequential swarm experiment injects faults progressively across stages. Each stage experiences a fault probability that increases with the stage index, simulating a deteriorating operational environment where mechanical wear and damage accumulate over time.

The fault probability for stage $i$ is defined as:
\begin{equation}
p_{\text{fault}}(i) = \min\left(p_0 + \left\lfloor \frac{i}{k} \right\rfloor \cdot \Delta p,\; p_{\text{max}}\right)
\end{equation}
where $p_0 = 0.01$ baseline probability, $\Delta p = 0.015$ increments, $k = 5$ tier bounds, and maximum probability saturates at 50\%. Mechanical damage scales with a zero-mean Gaussian tracking offset whose variance increases across matching operational tiers.

The progressive fault parameter configurations align with vectors mapped inside Table 6.1.

\begin{figure}[H]
  \centering
  \small
  \setlength{\unitlength}{1mm}
  \begin{picture}(140,50)
    % Blocks
    \put(10,30){\framebox(25,10){Agent $A_i$}}
    \put(55,30){\framebox(30,10){\texttt{factory\_state}}}
    \put(105,30){\framebox(25,10){Agent $A_{i+1}$}}
    % Labels
    \put(22,20){\framebox(40,8){Write Outcome}}
    \put(78,20){\framebox(40,8){Read History}}
    % Flow Lines
    \put(35,35){\vector(1,0){20}}
    \put(85,35){\vector(1,0){20}}
    \put(22,30){\line(0,-1){2}}
    \put(22,28){\vector(0,-1){0}}
    \put(117,30){\line(0,-1){2}}
  \end{picture}
  \caption{Temporal stigmergy execution block passing discrete environmental updates over time.}
  \label{fig:swarm_control_loop_native_ch6}
\end{figure}

\section{The Inter-Agent Communication Ledger as Stigmergic Channel}

\subsection{The Factory State Dictionary}

The dictionary layout and sequential read-modify-write tracking operations follow minimal simplicity boundaries, persistent memory allocations, and localized data readings.

\subsection{The Recovery Threshold Mechanism}

The trigger parameters dictate an active threshold value $\tau = 2$ consecutive failures. Recovery options override standard execution targets, shifting reaching coordinates to a stable central physical anchor, and relaxing tracking tolerances to mitigate geometric divergence bottlenecks.

\subsection{Comparison to Biological Stigmergy}

Algorithmic ledger states act as exact numerical functional analogs to biological chemical trace systems (pheromone deposits, sensory thresholds, decentralized reading nodes) compressing temporal footprints through digital permanence channels.

\section{State-Dependent Collective Adaptation: Experimental Results}

\subsection{Baseline: Standard IK Controller (No Stigmergy)}

Blind comparative loops execute single-robot chains independently, isolating reaching tasks from predecessors' histories. Disconnection from shared state pathways triggers cascading trial failures under progressive faults.

\subsection{NeuroKin with Temporal Stigmergy}

The active visual feedback loop reads diagnostic trace memory indices at initialization, enforcing localized parameter constraints to bypass error propagation paths cleanly across sequential nodes.

\subsection{Quantitative Comparison: 100-Stage Results}

Table~\ref{tab:swarm_results_ch6} summarizes the quantitative findings extracted across the 100-stage operational chain under progressive fault injections:
\begin{table}[H]
\centering
\caption{Sequential swarm performance comparison (100 stages, progressive fault injection).}
\label{tab:swarm_results_ch6}
\resizebox{\textwidth}{!}{%
\begin{tabular}{lcc}
\toprule
\textbf{Metric} & \textbf{Baseline (No Stigmergy)} & \textbf{NeuroKin (Temporal Stigmergy)} \\
\midrule
Overall success rate & 13\% & 88\% \\
Mean steps to convergence & 75.8 & 57.0 \\
Median steps to convergence & 48.5 & 34.5 \\
Failures (out of 100) & 87 & 12 \\
\bottomrule
\end{tabular}
}
\end{table}

\subsection{Analysis of Baseline Catastrophic Failure}

The baseline controller collapses to 13\% post-fault success once cascading parameters breach nominal limits, showing that decoupled structures fail to withstand progressive environmental degradation paths.

\subsection{Analysis of NeuroKin Stigmergic Adaptation}

NeuroKin channels preserve an 88\% success ceiling post-fault, confirming that real-time visual alignment overrides tracking divergence under active coordination mechanisms.

\subsection{Step Efficiency Analysis}

Figures~\ref{fig:swarm_comparison_ch6} and 4.10 from the validation script chart convergence step distributions. NeuroKin reduces tracking timeouts significantly, optimizing the median step footprint down by 29\%.

\begin{figure}[H]
    \centering
    \includegraphics[width=0.95\textwidth]{Neurokin vs Basline with fault.png}
    \caption{Algorithm Efficiency in Steps under injected faults. The IK baseline fails structurally (capped at 201 steps), while the NeuroKin swarm reads shared states and adapts behavior, maintaining efficiency between 20-40 steps.}
    \label{fig:swarm_comparison_ch6}
\end{figure}

\section{Statistical Analysis and Failure Mode Characterization}

Process survival analyses, operational breakdown metrics, and entropy channel information capacities track attributes detailed inside Sections 6.5.1, 6.5.2, and 6.5.3. Bounded ledger constraints isolate state dynamics with high information efficiency.

\section{Comparison to Prior Work on Swarm Coordination}

Coordination scalability patterns and reinforcement comparative dimensions align with structures discussed in Section 6.6.

\section{Limitations and Failure Case Analysis}

Concurrency control boundaries, ledger bottlenecks, fault profile limits, and recovery target locations follow parameters evaluated inside Section 6.7.

\section{Discussion: Validating the ``Toward Swarm Coordination'' Hypothesis}

Empirical capabilities prove the high relevance of sub-millisecond diagnosis towards establishing active coordination matrices. Bypassing processing latency limits prevents system state tracking lags, satisfying the computational prerequisite for unified multi-robot deployments.

\section{Summary}

This chapter proved that low-latency visual self-modeling provides the fundamental diagnostic backbone required for collective fault tolerance under cascading damage regimes. The next chapter introduces global evaluations.
\clearpage

% ==========================================
% CHAPTER 7: DISCUSSION
% ==========================================
\chapter{Discussion}
\label{chap:discussion_main_ch7}

This chapter critically evaluates the overarching hypothesis that motivated this thesis: that explicit 3D volumetric reconstruction is not a prerequisite for control-oriented self-modeling, and that sub-millisecond individual self-diagnosis enables decentralized swarm coordination via temporal stigmergy. We analyze the trade-offs inherent in the direct decoding paradigm, defend the sequential stigmergy approach against alternative coordination mechanisms, and discuss critical threats to validity---particularly the simulation-to-reality gap and the abstraction of mechanical degradation. Throughout, we ground the analysis in the empirical findings from Chapters 4--6, contextualizing them within the broader literature surveyed in Chapter 2.

\section{Trade-offs of Direct Decoding: Sacrificing Novel-View Synthesis for Extreme Inference Speed}

The central architectural decision in NeuroKin is the elimination of explicit 3D reconstruction. While this choice enabled the dramatic 1,418$\times$ speedup over FFKSM (7,400 FPS vs. 5.22 FPS) with a simultaneous +4.53 dB PSNR improvement, it carries inherent limitations that must be acknowledged and, where possible, mitigated.

\subsection{The Single-Viewpoint Constraint}

NeuroKin learns a fixed mapping from joint angles to the specific camera viewpoint used during training. Unlike volumetric methods that can support novel-view rendering, NeuroKin cannot generalize to cameras placed at positions not seen during training. This limitation is fundamental: the network has never been supervised on images from other angles, so it has no basis for predicting them.

In practice, many robotic applications tolerate this constraint. For collision checking with a fixed overhead camera, or for proprioceptive drift correction using the robot's own egocentric view, a single known viewpoint suffices \cite{hu2024teaching, chen2022full}. However, for tasks requiring multi-view awareness---such as grasping objects from arbitrary orientations or navigating through cluttered environments with moving cameras---the single-viewpoint limitation becomes restrictive.

Three potential mitigation strategies emerge from the literature. First, a multi-head architecture could train separate output heads for a discrete set of predefined viewpoints (front, side, top, etc.), amortizing the encoder cost across predictions while retaining the fast inference of direct decoding. Second, camera conditioning---concatenating camera parameters $\mathbf{c} \in \mathbb{R}^6$ (position and orientation unit vector) to the joint angle input---would enable the network to learn a conditional mapping $f(\theta, \mathbf{c}) \rightarrow I$, at the cost of increased training data and a modest latency penalty. Third, a hybrid approach could use NeuroKin for the primary viewpoint (highest frame rate) and sparse FFKSM evaluations for auxiliary views when needed, trading off speed for geometric completeness. Each of these directions is left for future work.

\subsection{Binary Silhouettes vs. RGB Prediction}

NeuroKin predicts binary masks rather than full RGB images. While silhouettes contain sufficient structural information for the control tasks evaluated in this thesis---collision checking, end-effector tracking, and damage detection---they discard texture, color, and material properties. For applications requiring visual inspection (e.g., detecting tool wear, surface damage, or identifying objects by color), the binary representation is insufficient.

The decision to use silhouettes was deliberate. Binary masks are robust to lighting variation when properly thresholded, require minimal preprocessing, and provide dense gradient supervision. Moreover, they focus the network on geometric structure rather than surface appearance, which is precisely what matters for kinematic self-modeling. ResNeuroKin-D's depth prediction head demonstrates that additional geometric cues can be learned as auxiliary tasks without sacrificing the core silhouette prediction speed. Extending this to RGB prediction is possible---the output head could be modified to produce three channels instead of one---but would increase the learning burden and likely reduce inference speed.

\subsection{The Pareto Frontier: Novel-View Synthesis vs. Control-Loop Latency}

Figure~\ref{fig:pareto_frontier_ch7} places NeuroKin, ResNeuroKin-D, K-3DGS, and FFKSM on a conceptual Pareto frontier with reconstruction quality (PSNR) on one axis and inference speed (FPS, logarithmic) on the other. The key insight from this analysis is that for applications requiring real-time control (shaded region above 1,000 FPS), direct decoding is the only viable architecture among those evaluated. In our comparison, the rendering-based baselines remain far below the 1,000 FPS region required for 1 kHz control \cite{hu2025learning}.

\begin{figure}[H]
\centering
\includegraphics[width=0.95\textwidth]{figures/pareto_frontier.png}
\caption{Pareto frontier illustrating the optimal trade-off between model accuracy (PSNR) and computational speed (FPS). NeuroKin occupies the upper-right region, achieving both superior quality and sub-millisecond latency. The red dashed line indicates the 1,000 FPS threshold for 1 kHz control loops.}
\label{fig:pareto_frontier_ch7}
\end{figure}

The practical implication is clear: for robots that must react to embodiment drift within control-loop deadlines (1--2 ms), volumetric reconstruction is not merely inefficient---it is fundamentally incompatible. Direct decoding, by contrast, provides state estimates with sufficient speed and fidelity to stabilize feedback control. This finding directly answers Research Question RQ1 and validates the central hypothesis that 3D reconstruction is not a prerequisite for control-oriented self-modeling. This architectural divergence becomes particularly clear when contrasted with recent efforts to integrate differentiable rendering directly into the control loop. Approaches that optimize trajectories by anchoring control directly to observation-space rendering \cite{liu2024differentiable} offer tight visual-motor coupling, yet they remain tethered to the fundamental computational bottleneck of the rendering pipeline. Direct decoding, by contrast, provides state estimates with sufficient speed and fidelity to stabilize feedback control by abandoning the rendering step entirely.

\subsection{Supervision Efficiency as the Key Explanatory Variable}

The 1,418$\times$ speedup with improved quality appears paradoxical: conventional wisdom suggests that speed and quality trade off. The resolution lies in supervision efficiency, as formalized in Section 4.1.2. NeuroKin's single forward pass provides dense gradient information across all 10,000 output pixels, whereas FFKSM's ray-marching provides each network evaluation with only 0.0156 pixels of supervision on average. This 640,000$\times$ difference in supervision density explains why NeuroKin trains 89$\times$ faster (33.5 seconds vs. 50 minutes) and generalizes better from limited data (2,000 samples vs. 12,000 in the original FFKSM paper \cite{hu2024teaching}).

Supervision efficiency is not merely an implementation detail; it is a fundamental property of the architectural choice. Any method that requires sparse sampling of a volumetric field (whether NeRF, 3DGS, or occupancy networks) will inherently suffer from low supervision density. Direct decoding, by contrast, achieves maximum parallelism and dense supervision by design. This insight suggests a general principle: when the downstream task requires only 2D predictions from a fixed viewpoint, 3D reconstruction is an unnecessary intermediate representation that imposes a severe computational tax.

\section{Validating the ``Toward Swarm Coordination'' Hypothesis: Temporal Stigmergy as a Prerequisite}

The thesis title includes the phrase ``Toward Swarm Coordination'' with deliberate scientific restraint. Chapter 6 demonstrated that a 100-stage sequential swarm, where agents read and write a shared failure ledger, achieves 88\% post-fault success compared to the baseline's 13\%. This section defends the claim that this temporal stigmergy mechanism validates the \textit{computational prerequisite} for physical swarm deployment, even though concurrent spatial coordination remains future work.

\subsection{Why Temporal Stigmergy Proves the Prerequisite}

A physical swarm of robots operating simultaneously in a shared workspace requires three foundational capabilities: (1) rapid individual self-diagnosis to detect faults and estimate state, (2) a mechanism for agents to share diagnostic information without centralized coordination, and (3) the ability to adapt behavior based on shared information to improve collective outcomes. Chapter 6 demonstrates all three capabilities in a controlled sequential setting:
\begin{enumerate}
    \item \textbf{Rapid self-diagnosis}: NeuroKin-3D provides end-effector position estimates with 0.400 ms latency and 2.39 mm mean error (from Section 4.4.2), enabling real-time failure detection.
    \item \textbf{Decentralized information sharing}: The \texttt{factory\_state} ledger serves as a stigmergic communication channel---agents write their outcomes and read predecessors' states without explicit peer-to-peer messaging.
    \item \textbf{Adaptive behavior}: The threshold rule (two consecutive failures triggers recovery mode) changes target selection from random workspace positions to a fixed conservative coordinate [0.15, 0.0, 0.15] and doubles the success tolerance from 1 cm to 2 cm, improving success rate from 13\% to 88\%.
\end{enumerate}

These capabilities are necessary conditions for spatial swarm coordination. If agents cannot diagnose themselves rapidly, they cannot provide reliable state estimates to the swarm. If they cannot read and write shared state, they cannot coordinate. If they cannot adapt based on history, they cannot achieve fault tolerance. By demonstrating these capabilities in a sequential pipeline, this thesis proves that the foundational building blocks of swarm intelligence are within reach for the NeuroKin architecture.

\subsection{The Role of Sub-Millisecond Inference in Stigmergic Coordination}

The 0.400 ms inference time of NeuroKin-3D is critical for the temporal stigmergy mechanism. Consider an alternative where self-diagnosis required 191.6 ms (FFKSM's latency from Chapter 3). In a 100-stage sequential pipeline, the cumulative diagnostic time alone would exceed 19 seconds, even before accounting for control and actuation. More fundamentally, an agent with 191.6 ms inference cannot reliably determine failure causes in real time---by the time the diagnosis completes, the system state may have changed.

Sub-millisecond inference enables three properties essential for stigmergy: (1) real-time failure detection (the agent knows within a single control cycle whether it has succeeded or failed); (2) immediate ledger updates (the \texttt{factory\_state} is updated before the next control cycle begins); and (3) rapid adaptation (subsequent agents read the updated state without perceptible delay). Without sub-millisecond inference, temporal stigmergy would introduce latency proportional to the number of agents, causing the coordination mechanism to become the bottleneck.

\subsection{Comparison to Alternative Coordination Mechanisms}

Table 7.1 situates temporal stigmergy relative to other swarm coordination paradigms discussed in the literature \cite{bongard2006resilient, cully2015robots, domberg2026selfadapting}.

Our temporal stigmergy approach preserves the low-complexity spirit of biological stigmergy while operating at digital time scales. Unlike pheromone-based stigmergy, which requires physical trace persistence and decay \cite{bongard2006resilient}, temporal stigmergy uses persistent digital memory. Unlike richer coordination schemes, the ledger relies on a deliberately small shared state and simple update rule.

Critically, temporal stigmergy requires no training beyond the individual NeuroKin self-model; the coordination mechanism is a simple threshold rule operating on the shared ledger. This simplicity is a feature: the decision to enter recovery mode is transparent and verifiable, and the mechanism works with any agent that can read and write the ledger.

\subsection{Limitations of the Sequential Paradigm}

The most significant limitation of Chapter 6 is that temporal stigmergy coordinates agents across time, not space. The reference implementation executes agents sequentially, not concurrently. Concurrent swarm deployment introduces challenges not addressed here: collision avoidance, state sharing, mapping, and distributed task allocation under resource contention \cite{matsuki2024gaussian, keetha2024splattam}. These are separate research problems beyond this thesis's scope.

Temporal stigmergy is positioned as a \textit{computational prerequisite} for swarm coordination---demonstrating that fast individual self-diagnosis can support emergent coordination---not as a complete solution for concurrent swarm deployment. The transition from sequential to concurrent deployment, while substantial, is an engineering challenge rather than a fundamental research question. Future work must address concurrency control (atomic ledger operations), communication latency (non-instantaneous updates), and spatial coordination. 

To achieve true spatial stigmergy where multiple robots move together safely, the swarm will require reliable decentralized mapping, collision avoidance, and shared world-modeling. Those ingredients remain future work outside the present sequential validation. Recent advancements in feature-enriched 3D Gaussian splatting SLAM \cite{thirgood2025featureslam} and uncertainty-aware RGB-D SLAM \cite{tran2025varsplat} provide the exact low-latency, semantic frameworks necessary to handle multi-agent tracking, collision avoidance, and shared world-modeling. Furthermore, persistent spatial memory using 3D Gaussian splatting, such as GSMem \cite{lu2025gsmem}, offers a promising direction for maintaining re-renderable scene representations that multiple agents could concurrently query for coordination tasks.

\subsection{Information Bottleneck and Recovery Mode Design}

The shared ledger contains only the \texttt{consecutive\_failures} counter, providing at most $\log_2(3) \approx 1.58$ bits of information. This minimal representation proved sufficient for the reaching task under progressive fault injection, achieving 88\% success. However, for more complex coordination scenarios---task allocation, resource tracking, formation maintenance---richer state information (fault type, severity, location) may be required. Extending the ledger to support multi-dimensional state is a direction for future work.

The recovery mode implemented in Chapter 6 uses a fixed conservative target [0.15, 0.0, 0.15] and doubles the success tolerance from 1 cm to 2 cm. This design, while effective for the 4-DOF arm, may be suboptimal for other morphologies or tasks. Adaptive recovery target selection based on estimated fault parameters (e.g., using the joint angle predictions from NeuroKin-3D to infer which joint is damaged) could yield faster recovery. The 12\% failure cases in Chapter 5---extreme fault magnitudes (>25$^\circ$), workspace exit, and silhouette degradation---suggest specific improvements: anti-windup for the integral term, a retreat-and-reapproach planner, and multi-camera redundancy or occlusion-aware training.

\section{Threats to Validity and Directions for Mitigation}

\subsection{Simulation-to-Reality Gap}

All experiments in this thesis were conducted in PyBullet simulation with perfect binary silhouettes (thresholded at 240). Real-world deployment introduces challenges that systematically threaten the validity of our conclusions: lighting variation, background clutter, camera noise, and calibration errors. The literature on sim-to-real transfer \cite{tan2018simtoreal, aljalbout2026realitygap} identifies four distinct gap dimensions: dynamics mismatch (physics simulation divergence), perception mismatch (sensor modeling inaccuracy), actuation mismatch (motor characteristics divergence), and system construction mismatch (morphological/structural misalignment).

For NeuroKin, the most critical gap is perception. The network was trained on synthetic binary masks derived from perfect RGB rendering. Physical cameras produce images with shadows, specular highlights, motion blur, and sensor noise. Background subtraction or learned segmentation would be required to extract clean silhouettes in real environments. Domain randomization during training can improve robustness.

More broadly, bridging the visual perception gap will require stronger sim-to-real pipelines, tighter calibration, and deployment-time correction when the twin drifts from hardware behavior. SplatSim \cite{qureshi2025splatsim} and semantic 3D Gaussian splatting \cite{tang2025semantic} illustrate the direction of travel, but full validation on physical hardware remains necessary before claiming deployment readiness.

The stereoscopic setup also requires accurate extrinsic calibration. Calibration errors of 1-2 mm are acceptable given NeuroKin-3D's 2.39 mm mean spatial error, but errors exceeding 5 mm would degrade performance. Standard checkerboard calibration can achieve sub-millimeter accuracy, but the cameras must remain fixed relative to the robot base.

\subsection{Abstraction of Mechanical Degradation and Rigid-Body Assumptions}

The fault model used in Chapters 5 and 6---encoder slip on joint 2 with zero-mean Gaussian magnitude---represents a specific, relatively benign form of embodiment drift. Real hardware failures include complete joint seizure (zero degrees of freedom), broken gear teeth (periodic position errors), thermal drift (nonlinear, temperature-dependent), and communication packet loss (stochastic, bursty). The temporal stigmergy mechanism should generalize to these fault types, as it responds to task-level success/failure rather than specific fault signatures. However, the recovery mode's effectiveness may degrade if faults produce systematic biases that the visual estimator cannot track (e.g., a bent link that changes the kinematic tree).

The 12\% failure cases in Chapter 5 provide a taxonomy of scenarios where visual estimation alone is insufficient: extreme fault magnitudes (>25$^\circ$), workspace exit, and silhouette degradation. These failure modes are not unique to NeuroKin; they would affect any vision-based controller. Addressing them requires complementary mechanisms: anti-windup integrators, recovery planners, and redundant sensing.

It is also critical to explicitly acknowledge the boundary of the rigid-body assumption inherent in the current NeuroKin formulation. Direct silhouette mapping works exceptionally well for the rigid kinematic links and joint anomalies of the 4-DOF arm. However, extending visual self-modeling to soft robotics or embodiments with flexible elements would violate these rigid geometric assumptions. For such applications, physics-informed dynamic surrogates would be necessary to capture non-rigid morphological drift. Extending visual self-modeling to continuum or soft robotics requires representations that track the continuous dynamics of deformation rather than discrete geometric snapshots.

\subsection{Generalization to Higher Degrees of Freedom}

The 4-DOF serial manipulator used throughout this thesis is a simplified testbed. Scaling to high-DOF systems (e.g., humanoid robots with 30+ joints) introduces curse-of-dimensionality challenges. For instance, successfully modeling the expanded configuration space of 7-DOF manipulators (e.g., Franka Panda) has been shown to require high-DOF dynamic neural fields coupled with curricular data sampling \cite{schulze2024high}. The 1024-dimensional latent embedding may be insufficient for 30-dimensional joint spaces; covering high-dimensional workspaces requires exponentially more training samples; and self-occlusion becomes severe, requiring multi-camera fusion. Potential architectural extensions include hierarchical encoders (separate networks for limbs, torso, head), attention mechanisms to handle variable-length kinematic chains, and graph neural networks that exploit the robot's kinematic tree structure. Dynamics-grounded priors, such as the articulated-body dynamics network, offer a promising direction for sample-efficient scaling by embedding physical structure directly into the learning architecture. These extensions are beyond the scope of this thesis but represent important future work.

\subsection{Catastrophic Forgetting Under Structural Damage}

While NeuroKin-3D achieved 2.39 mm mean error under nominal conditions and recovered from encoder slip via visual feedback, structural damage (bent link, seized joint) required fine-tuning on 100 post-damage samples (5 seconds) to restore accuracy. Preliminary experiments with naive fine-tuning revealed catastrophic forgetting: updating the network on damaged configurations degraded predictions for healthy configurations. This is the stability-plasticity dilemma discussed in Section 2.2.3.

Mitigation strategies include Elastic Weight Consolidation (EWC), which protects important weights via Fisher information; sparse subnetworks identified via the Lottery Ticket Hypothesis; and mixture-of-experts architectures that isolate damage-specific modules. Each approach has trade-offs: EWC is computationally expensive for real-time adaptation, sparse subnetworks require careful initialization, and mixture-of-experts increases parameter count. For the 33.5-second training time of NeuroKin, a simpler strategy---periodic full retraining on a rolling window of recent observations---may be practical for many applications.

\subsection{Data Efficiency and Limited Training Samples}

All comparative experiments trained on 2,000 samples---one-sixth the scale of the original FFKSM paper's 12,000 samples. While this enabled fair controlled comparison, it raises questions about generalization to larger datasets. Our data efficiency analysis suggests NeuroKin benefits more from limited data than volumetric rendering: it achieved 21.88 dB with 2,000 samples, approaching FFKSM's reported 23+ dB quality despite 6$\times$ less data. Extrapolating to 12,000 samples, we hypothesize NeuroKin would achieve 25--26 dB, exceeding FFKSM even at full scale. However, this extrapolation remains untested and should be validated in future work.

\section{Summary of Discussion}

This chapter critically evaluated the trade-offs and validity of the contributions presented in Chapters 4--6. The key insights are:
\begin{enumerate}
    \item \textbf{Direct decoding trades novel-view synthesis for control-loop latency}. For applications requiring real-time response from a fixed viewpoint, this trade-off is highly favorable. For multi-view tasks, hybrid architectures or camera conditioning are necessary.
    \item \textbf{Supervision efficiency explains the paradoxical speed-quality improvement}. Dense gradients from direct pixel prediction provide 640,000$\times$ more supervision per network evaluation than volumetric ray-marching.
    \item \textbf{Temporal stigmergy validates the computational prerequisite for swarm coordination}. Sub-millisecond individual self-diagnosis, decentralized ledger updates, and threshold-based adaptation collectively enable 88\% post-fault success in a 100-stage sequential pipeline. Extending this coordination framework to concurrent multi-agent settings will require reliable shared world models and spatial coordination mechanisms.
    \item \textbf{Sim-to-real transfer remains the primary threat to validity}. Domain randomization, robust silhouette extraction, tighter calibration, and hardware validation are essential next steps. Zero-shot sim2real methods illustration promises directions, but deployment evidence remains decisive.
    \item \textbf{Catastrophic forgetting under structural damage is partially mitigated by rapid fine-tuning}, but principled continual learning methods are needed for long-term deployment. Furthermore, while the current architecture is well-suited for rigid bodies, scaling to soft robotics will necessitate physics-informed dynamic surrogates.
    \item \textbf{Scaling to higher degrees of freedom} (such as 7-DOF arms \cite{schulze2024high}) will require hierarchical encoders, stronger structural priors, and more effective handling of self-occlusion and configuration-space complexity.
\end{enumerate}

The following chapter synthesizes these findings into concrete conclusions and maps a technical roadmap for transitioning from sequential validation to physical swarm deployment.
\clearpage

% ==========================================
% CHAPTER 8: CONCLUSION AND FUTURE WORK
% ==========================================
\chapter{Conclusion and Future Work}
\label{chap:conclusion_future_work_ch8}

This chapter synthesizes the principal findings of this thesis, directly answering the four research questions posed in Chapter 1, and outlines a concrete technical roadmap for transitioning the validated sequential capabilities into physically embodied, spatially concurrent swarm deployments.

\section{Summary of Contributions and Research Findings}

This thesis addressed the central problem that explicit 3D volumetric reconstruction, while enabling high-fidelity visual self-modeling, imposes a latency barrier that fundamentally precludes real-time control integration. The core hypothesis---that sub-millisecond visual self-modeling can be achieved by bypassing explicit 3D reconstruction in favor of direct sensorimotor decoding, and that this capability enables decentralized swarm coordination via temporal stigmergy---has been rigorously validated through controlled simulation experiments. The findings are structured around the four research questions introduced in Chapter 1.

\subsection{RQ1: Architectural Efficiency --- Direct Decoding Defeats the Latency Paradox}

\textbf{Question}: To what extent can the computational latency of visual robot self-modeling be reduced by replacing 3D volumetric rendering with 2D direct sensorimotor decoding?

\textbf{Answer}: Direct sensorimotor decoding reduces latency by over three orders of magnitude while simultaneously improving reconstruction quality. The empirical evidence from Chapters 3 and 4 demonstrates that NeuroKin achieves 7,400 FPS (0.135 ms latency) and 21.88 dB PSNR---a 1,418$\times$ speedup over the implicit NeRF-based FFKSM (5.22 FPS, 17.35 dB PSNR) and a 200$\times$ speedup over the explicit 3DGS-based K-3DGS (37 FPS, 17.01 dB PSNR).

The resolution of the apparent speed-quality paradox lies in supervision efficiency. NeuroKin's single forward pass provides dense gradient information across 10,000 pixels simultaneously, whereas FFKSM's ray-marching provides each network evaluation with only 0.0156 pixels of supervision on average. This 640,000$\times$ difference in supervision density explains why NeuroKin trains 89$\times$ faster (33.5 seconds vs. 50 minutes) and generalizes more effectively from limited data (2,000 samples vs. the 12,000 samples required by the original FFKSM paper \cite{hu2024teaching}).

This finding directly challenges the conventional wisdom that 3D reconstruction is a necessary prerequisite for visual self-modeling. For control-oriented tasks requiring only silhouette-level awareness---collision checking, proprioceptive drift correction, and damage detection---the representational overhead of volumetric methods is computational waste. The principle of minimal representation emerges as a design guideline: neural self-models should maintain only the complexity required by downstream tasks.

\subsection{RQ2: Control Integration --- Sub-Millisecond Visual Servoing Enables Dynamic Stabilization}

\textbf{Question}: Can a visual self-model operating purely on direct convolutional decoding provide spatial estimates with sufficient accuracy and latency to actively stabilize a Jacobian-based inverse kinematics control loop?

\textbf{Answer}: Yes. Chapter 5 demonstrated that integrating NeuroKin-3D's stereoscopic spatial estimates (2.39 mm mean error, 0.400 ms latency) into a resolved-rate feedback loop enables real-time closed-loop control. Under nominal conditions, the visual controller converged 15\% faster than the proprioceptive baseline (35.9 vs. 42.2 steps, $p = 0.003$) with 18\% lower final error (2.3 vs. 2.8 mm, $p = 0.017$). This improvement stems from two factors: direct Cartesian measurement, which bypasses the accumulation of forward kinematic errors, and exponential smoothing ($\alpha = 0.6$), which attenuates high-frequency noise.

Critically, when a severe encoder slip fault was injected at step 50 (simulating calibration drift or thermal deformation), the proprioceptive baseline collapsed to 13\% success, while the NeuroKin visual controller maintained 88\% success. The recovery was rapid: the controller re-converged to sub-centimeter accuracy within 38.5 steps on average (0.77 seconds). The 12\% failure rate was dominated by extreme fault magnitudes ($>25^\circ$) causing joint limit saturation and workspace exit conditions, not by visual estimation failure.

A Lyapunov stability analysis confirmed that the proportional-only visual servo is exponentially stable. The estimation error $\boldsymbol{\delta}$ (bounded by 2.39 mm) enters as an additive disturbance, and the system converges to a residual error ball whose radius is proportional to $\|\boldsymbol{\delta}\|/K_p$. With $K_p = 0.25$, the predicted steady-state error of approximately 9.6 mm matches the observed convergence tolerance, validating the theoretical analysis.

\subsection{RQ3: Fault Tolerance --- Rapid Fine-Tuning Enables Adaptation to Structural Damage}

\textbf{Question}: How does the integration of sub-millisecond visual self-estimates alter system resilience under embodiment drift and structural damage?

\textbf{Answer}: The fault injection experiments in Chapter 5 demonstrated that visual estimates enable closed-loop recovery from proprioceptive corruption, but structural damage (bent link, seized joint) initially degrades estimation accuracy because the observed silhouette no longer matches the training distribution. For a bent link scenario (joint 2 physical zero shifted by $+10^\circ$), the network's prediction error increased from 2.4 mm to 12.7 mm mean error. For a seized joint scenario (joint 3 locked), the error diverged to 18.4 mm.

However, NeuroKin-3D's efficient training (33.5 seconds from scratch) enables rapid online adaptation. Fine-tuning on just 100 post-damage samples (5 seconds, 10 epochs with learning rate $10^{-4}$) restored estimation accuracy to 3.1 mm for the bent link scenario and 4.2 mm for the seized joint scenario. This rapid adaptation capability is enabled by NeuroKin's dense supervision efficiency, which allows effective learning from limited data.

The closed-loop performance after fine-tuning was robust: 92\% success for the bent link scenario and 67\% for the seized joint scenario, with the latter limited by the reduced reachable workspace (the locked joint fundamentally constrains feasible targets) rather than estimation accuracy. Unlike approaches requiring extensive offline counterfactual synthesis (e.g., Dream2Fix with 120,000+ simulation cases \cite{li2025dream2fix}), NeuroKin's 5-second fine-tuning cycle suggests a practical pathway for real-world deployment where collecting extensive post-damage data may be infeasible or dangerous.

\subsection{RQ4: Collective Adaptation --- Temporal Stigmergy Validates the Prerequisite for Swarm Coordination}

\textbf{Question}: Can the individual computational efficiency of a fast self-model enable decentralized, state-dependent adaptation across a sequential pipeline of agents?

\textbf{Answer}: Yes. Chapter 6 demonstrated that a 100-stage sequential swarm, where agents read and write a shared \texttt{factory\_state} ledger tracking consecutive failures, achieves 88\% post-fault success compared to the baseline's 13\%. The temporal stigmergy mechanism---agents coordinate across time by leaving environmental traces (failure counters) that trigger adaptive responses in subsequent agents---requires no centralized coordination, no explicit inter-agent communication, and no additional training beyond the individual NeuroKin self-model.

The threshold rule ($\tau = 2$ consecutive failures triggers recovery mode) is simple yet effective. A single failure ($\tau = 1$) would trigger recovery mode too aggressively, causing unnecessary conservative behavior after isolated, unrepeated faults. A higher threshold ($\tau = 3$ or $\tau = 4$) would delay response, allowing cascading failures to propagate. The empirically calibrated $\tau = 2$ balances these competing pressures.

The information efficiency is striking: with at most $\log_2(3) \approx 1.58$ bits of shared information (the 3-state ledger), the NeuroKin swarm achieves a 75 percentage point improvement over the baseline. The median steps to convergence for NeuroKin agents was 34.5, compared to 48.5 for the baseline's successful runs (a 29\% reduction), demonstrating that the recovery mode not only improves success rates but also operational efficiency.

Crucially, the 0.400 ms inference time of NeuroKin-3D is essential for this mechanism. It enables real-time failure detection within a single control cycle, immediate ledger updates before the next agent is instantiated, and rapid adaptation without stale state windows. Without sub-millisecond inference, temporal stigmergy would introduce latency proportional to the number of agents, causing the coordination mechanism to become the bottleneck.

\section{Technical Contributions Summary}

Beyond answering the research questions, this thesis makes four specific technical contributions to the fields of robotic self-modeling, neural representation, and decentralized multi-agent systems:

\subsection{Contribution 1: Formal Benchmarking of the Latency Paradox}

Through rigorous ground-up reproduction and quantitative evaluation of both implicit (FFKSM/NeRF) and explicit (K-3DGS) neural rendering self-models on a standardized 2,000-sample dataset, this research formally exposes their control incompatibility. The exact latency metrics---191.6 ms for FFKSM, 27.0 ms for K-3DGS---establish the mathematical limits that necessitate a paradigm shift. The reproduction also uncovered three undocumented implementation challenges (black-screen convergence, kinematic blindness, and weighted loss tuning) that are essential for successful training but omitted from the original publication \cite{hu2024teaching}.

\subsection{Contribution 2: The NeuroKin Direct Sensorimotor Decoder}

NeuroKin is a lightweight, fully convolutional direct sensorimotor decoder that completely bypasses volumetric reconstruction. Through extensive empirical validation, NeuroKin achieves 21.88 dB PSNR at 7,400 FPS (0.135 ms latency)---a 1,418$\times$ speedup over implicit NeRF baselines with a +4.53 dB quality improvement. The stereoscopic extension, NeuroKin-3D, recovers 3D spatial coordinates with 2.39 mm mean error at 2,500 FPS (0.400 ms latency). The multi-task variant, ResNeuroKin-D, adds depth prediction (21.23 dB PSNR at 2,500 FPS) and induces geometric structure in the latent representation. All architectures are open-sourced for reproducibility.

\subsection{Contribution 3: Visual Closed-Loop Damage Adaptation}

This research successfully synthesizes NeuroKin's sub-millisecond spatial estimates with a Jacobian-based inverse-kinematics feedback loop. The integration enables dynamic error recovery after severe fault injection (88\% success vs. 13\% baseline). Lyapunov stability analysis confirms exponential stability, and the empirical demonstration that high-speed visual estimates can override corrupted proprioceptive data proves that sub-millisecond self-modeling enables real-time closed-loop adaptation.

\subsection{Contribution 4: Temporal Stigmergy for Sequential Swarm Coordination}

The validation of stigmergic swarm coordination under compounding mechanical faults demonstrates that ultra-fast individual self-diagnosis is the computational prerequisite for swarm intelligence. The 100-stage sequential swarm with progressive fault injection (fault probability from 1\% to 29.5\%, fault scale from 0.005 rad to 0.062 rad) achieves 88\% post-fault success through a decentralized ledger with only 1.58 bits of shared information. The mechanism requires no centralized coordination, no explicit communication, and no additional training beyond the individual self-model.

\section{Future Roadmap: From Sequential Validation to Spatial Swarms}

The thesis title specifies ``Toward Swarm Coordination'' deliberately. While Chapter 6 validated the computational prerequisites for swarm intelligence in a sequential pipeline, transitioning to concurrent spatial multi-agent deployment requires addressing substantial additional challenges. This section maps a concrete technical roadmap for future work, organized by time horizon.

\subsection{Immediate Next Steps (6--12 Months): Bridging the Sim-to-Real Gap}

The most urgent priority is hardware validation. All experiments in this thesis were conducted in PyBullet simulation with perfect binary silhouettes (thresholded at 240). Real-world deployment introduces challenges that must be systematically addressed:

\paragraph{Robust Silhouette Extraction}
Threshold-based segmentation fails under non-uniform lighting and background clutter. Future work should train a segmentation network on a dataset of real robot images with diverse backgrounds, lighting conditions, and occlusions. Domain randomization during training---randomly varying brightness, contrast, adding synthetic shadows, and applying random occlusions---can improve robustness. The exponential smoothing filter should be retained to attenuate camera noise and segmentation artifacts.

\paragraph{Stereoscopic Calibration}
The dual-camera system requires accurate extrinsic calibration (relative position and orientation of the two cameras). Calibration errors of 1-2 mm are acceptable given NeuroKin-3D's 2.39 mm mean error, but errors exceeding 5 mm would degrade performance. Standard checkerboard calibration can achieve sub-millimeter accuracy. For mobile platforms where cameras may move relative to the robot base, online calibration methods will be required.

\paragraph{Embedded Deployment}
The Tesla T4 GPU used in simulation provides 8.1 TFLOPS of FP32 compute. Embedded deployment on edge devices like NVIDIA Jetson Orin (5.3 TFLOPS, 15-25W) may reduce inference speed but should still exceed 1,000 FPS given NeuroKin's efficiency (7,400 FPS on T4). Future work should benchmark inference latency on Orin and implement quantization (FP16 or INT8) to further reduce latency and power consumption. The Jetson Orin's memory bandwidth constraints (204 GB/s) are less restrictive than the T4's (320 GB/s), but the parallel radix sort operations required for some vision pipelines are memory-bandwidth bound; NeuroKin's convolutional architecture avoids these bottlenecks.

\paragraph{Continual Learning for Damage Adaptation}
While fine-tuning on 100 post-damage samples restored accuracy in simulation, real hardware requires principled continual learning methods to prevent catastrophic forgetting. Future work should implement and evaluate Elastic Weight Consolidation (EWC), which protects important weights via the Fisher information matrix; sparse subnetwork methods based on the Lottery Ticket Hypothesis, which isolate damage-specific updates; and mixture-of-experts architectures, which maintain separate subnetworks for healthy and damaged configurations. For NeuroKin's 33.5-second training time, periodic full retraining on a rolling window of recent observations may also be practical.

\subsection{Medium-Term Extensions (1--2 Years): Concurrent Spatial Swarms}

With hardware validation complete, the next phase is transitioning from temporal (sequential) stigmergy to spatial (concurrent) stigmergy:

\paragraph{Concurrency Control for the Shared Ledger}
The sequential pipeline used a single \texttt{factory\_state} variable updated by one agent at a time. Concurrent swarms require atomic operations or consensus mechanisms to prevent race conditions. Future work should implement a distributed ledger (e.g., using Redis or a lightweight consensus protocol such as Raft) that supports atomic read-modify-write operations. For bandwidth-constrained environments, the ledger could be gossip-based with eventual consistency, which is sufficient for fault tolerance as the system can tolerate stale state (a few control cycles) without catastrophic failure.

\paragraph{Decentralized SLAM and Spatial Coordination}
Concurrent agents must avoid collisions while coordinating via shared state. Future work should integrate NeuroKin's visual self-modeling with decentralized mapping and a shared world model. Collision avoidance can be implemented as a simple rule: agents with higher consecutive failure counts yield to agents with lower counts, or agents in recovery mode move to safe coordinates. This preserves the stigmergic principle of indirect coordination without centralized planning.

\paragraph{Scalability Analysis}
The lightweight state-sharing pattern of temporal stigmergy should extend naturally to spatial swarms if the ledger is sharded by region or task. Future work should empirically evaluate success rate and communication overhead for swarms of 10, 50, and 100 agents operating concurrently in a shared workspace.

\paragraph{Fault Model Generalization}
The encoder slip fault model used in this thesis is a specific, relatively benign form of embodiment drift. Real hardware failures include complete joint seizure (zero degrees of freedom), broken gear teeth (periodic position errors), thermal drift (nonlinear, temperature-dependent), and communication packet loss (stochastic, bursty). The temporal stigmergy mechanism should generalize to these fault types as it responds to task-level success/failure rather than specific fault signatures. However, the recovery mode's effectiveness may degrade if faults produce systematic biases that the visual estimator cannot track (e.g., a bent link that changes the kinematic tree). Future work should evaluate the mechanism under a broader fault model.

\subsection{Long-Term Vision (2--5 Years): Adaptive Swarm Intelligence}

The ultimate goal is a swarm of robots that collectively adapts to embodiment drift, damage, and environmental changes without centralized control:

\paragraph{Hierarchical Stigmergy}
The current mechanism uses a single global ledger. Future work should explore hierarchical stigmergy: local ledgers for within-region coordination and a global ledger for cross-region coordination. This mirrors biological stigmergy where local pheromone trails aggregate into global patterns. The hierarchy would enable swarms to scale to hundreds of agents while maintaining communication efficiency.

\paragraph{Learned Recovery Policies}
The current recovery mode uses a fixed conservative target and doubles the success tolerance from 1 cm to 2 cm. Future work should learn adaptive recovery policies using reinforcement learning, where the policy conditions on the shared ledger state and the robot's estimated fault parameters (from NeuroKin-3D's joint angle predictions). The policy could learn to select optimal recovery targets, tolerance levels, and even exploration strategies for fault diagnosis. The 33.5-second training time of NeuroKin enables rapid policy iteration.

\paragraph{Cross-Embodiment Transfer}
Future work should evaluate whether NeuroKin models trained on one robot morphology can transfer to another via fine-tuning on limited data. The data efficiency advantage (21.88 dB with 2,000 samples) suggests that few-shot transfer is feasible. This would enable a swarm of heterogeneous robots to share self-modeling knowledge. Architectures that encode kinematic connectivity rather than visual appearance alone could further improve transfer performance.

\paragraph{Human-Swarm Interaction}
The interpretability of the stigmergic ledger (\texttt{consecutive\_failures} counter) enables natural human-swarm interaction. Future work should explore interfaces where human operators can inspect the ledger, manually reset counters, or inject recovery mode commands. This would provide a safety backup for autonomous swarm operations. The ledger state is compact and the threshold rule is transparent and verifiable.

\section{Future Roadmap: From Sequential Validation to Spatial Swarms}

The future roadmap is strictly bounded by hardware realities. Because physical deployment was impossible given the temporal constraints of this Master's degree, the next step is rigorous hardware validation under realistic sensing and calibration conditions. This foundational bridging of the sim-to-real gap will serve as the immediate launching point for the author's forthcoming doctoral research. Following that, concurrent multi-robot experiments that add collision avoidance and atomic shared-state updates can be effectively prototyped. Those extensions are consistent with the results in this thesis, but they are not claimed as solved here.

\section{Concluding Remarks}

This thesis demonstrates that fast visual self-modeling is useful in practice: it supports real-time closed-loop correction when proprioception is corrupted and improves a sequential multi-agent pipeline through a minimal shared failure ledger. The results do not claim full spatial swarm deployment, but they do establish the computational foundation needed for it.

The central design lesson is to keep the representation as small as the task allows. For the control and coordination problems studied here, direct decoding is sufficient, volumetric reconstruction is too slow, and a minimal ledger is enough to support temporal stigmergy. Those are the concrete conclusions supported by the experiments in this thesis.
\clearpage

\begin{thebibliography}{999}

\bibitem{bongard2006resilient}

J. Bongard, V. Zykov, and H. Lipson, ``Resilient machines through continuous self-modeling,'' \textit{Science}, vol. 314, no. 5802, pp. 1118--1121, 2006.

\bibitem{cully2015robots}

A. Cully, J. Clune, D. Tarapore, and J.-B. Mouret, ``Robots that can adapt like animals,'' \textit{Nature}, vol. 521, no. 7553, pp. 503--507, 2015.

\bibitem{kwiatkowski2019task}

R. Kwiatkowski and H. Lipson, ``Task-agnostic self-modeling machines,'' \textit{Science Robotics}, vol. 4, no. 26, p. eaau9354, 2019.

\bibitem{chen2022full}

B. Chen, R. Kwiatkowski, C. Vondrick, and H. Lipson, ``Full-body visual self-modeling of robot morphologies,'' \textit{Science Robotics}, vol. 7, no. 68, p. eabn1944, 2022.

\bibitem{hu2024teaching}

Y. Hu, J. Lin, and H. Lipson, ``Teaching robots to build simulations of themselves,'' \textit{Nature Machine Intelligence}, vol. 6, pp. 123--130, 2024.

\bibitem{hu2025learning}

K. Hu, P. Yu, and N. Tan, ``Learning high-fidelity robot self-model with articulated {3D} Gaussian splatting,'' \textit{Int. J. Robotics Research}, 2025.

\bibitem{tan2018simtoreal}

J. Tan \textit{et al.}, ``{Sim-to-Real}: Learning agile locomotion for quadruped robots,'' in \textit{Proc. Robot.: Sci. Syst. (RSS)}, 2018.

\bibitem{heiden2022inferring}

E. Heiden, Z. Liu, V. Vineet, E. Coumans, and G. S. Sukhatme, ``Inferring articulated rigid body dynamics from RGBD video,'' in \textit{Proc. IEEE Int. Conf. Robot. Autom. (ICRA)}, 2023.

\bibitem{lou2026drex}

H. Lou \textit{et al.}, ``D-REX: Differentiable real-to-sim-to-real engine for learning dexterous grasping,'' in \textit{Proc. Int. Conf. Learn. Represent. (ICLR)}, 2026.

\bibitem{qureshi2025splatsim}

M. N. Qureshi, Y. Tang, Z. Zhang, X. Li, Y. Wang, and A. Gupta, ``{SplatSim}: Zero-shot sim2real transfer of {RGB} manipulation policies using Gaussian splatting,'' arXiv preprint arXiv:2404.13099, 2025.

\bibitem{tang2025semantic}

J. Tang, Z. He, K. Zhang, Y. Liu, M. Chen, and W. Xu, ``Bridging simulation and reality: Cross-domain transfer with semantic 2D Gaussian splatting,'' arXiv preprint arXiv:2512.04731, 2025.

\bibitem{yang2025robosplat}

S. Yang, K. Zhan, Y. Zhang, L. Lin, X. Huang, and M. Li, ``Novel demonstration generation with Gaussian splatting enables robust one-shot manipulation,'' arXiv preprint arXiv:2504.13175, 2025.

\bibitem{yu2025real2render2real}

J. Yu, H. Su, D. Fridovich-Keil, and A. Gupta, ``{Real2Render2Real}: Scaling robot data without dynamics simulation or robot hardware,'' arXiv preprint arXiv:2505.09601, 2025.

\bibitem{abou2024realissim}

J. Abou-Chakra, Y. Zhu, L. Zhang, D. Fridovich-Keil, and P. Agrawal, ``Real-is-sim: Bridging the sim-to-real gap with a dynamic digital twin,'' arXiv preprint arXiv:2504.03597, 2025.

\bibitem{jain2025polaris}

A. Jain, R. Sharma, S. Kumar, S. Kolathaya, Z. Zhang, and D. Song, ``{PolaRiS}: Scalable real-to-sim evaluations for generalist robot policies,'' arXiv preprint arXiv:2512.16881, 2025.

\bibitem{li2024simpler}

X. Li \textit{et al.}, ``Evaluating real-world robot manipulation policies in simulation,'' in \textit{Proc. Conf. Robot. Learn. (CoRL)}, 2024.

\bibitem{grambow2026anomaly}

N. Grambow \textit{et al.}, ``Anomaly detection for generic failure monitoring in robotic assembly, screwing and manipulation,'' \textit{IEEE Robotics and Automation Letters}, vol. 11, no. 5, pp. 3664--3671, 2026.

\bibitem{fu2025contrastive}

Y. Fu, X. Chen, B. Zhang, M. Yang, and K. Zhang, ``Contrastive forward prediction reinforcement learning for adaptive fault-tolerant legged robots,'' in \textit{Proc. Conf. Robot. Learn. (CoRL)}, 2025.

\bibitem{chen2024vlm}

H. Chen, Y. Yao, R. Liu, C. Liu, and J. Ichnowski, ``Automating robot failure recovery using vision-language models with optimized prompts,'' arXiv preprint arXiv:2409.03966, 2024.

\bibitem{li2025dream2fix}

D. Li \textit{et al.}, ``Learning actionable manipulation recovery via counterfactual failure synthesis,'' arXiv preprint arXiv:2603.13528, 2026.

\bibitem{briscoe2025deft}

G. G. Briscoe-Martinez, Y. Gautam, R. Shetty, A. Pasricha, M. M. Nicotra, and A. Roncone, ``Moving on, even when you're broken: Fail-active trajectory generation via diffusion policies conditioned on embodiment and task,'' arXiv preprint arXiv:2602.02895, 2026.

\bibitem{poddar2026embedding}

N. Poddar, S. McCrory, L. Penco, G. Clark, H. E. Sevil, and R. Griffin, ``Embedding classical balance control principles in reinforcement learning for humanoid recovery,'' arXiv preprint arXiv:2603.08619, 2026.

\bibitem{li2026pchc}

H. Li \textit{et al.}, ``{PCHC}: Enabling preference-conditioned humanoid control via multi-objective reinforcement learning,'' arXiv preprint arXiv:2603.24047, 2026.

\bibitem{lee2026tolebi}

H. Lee, W.-J. Baek, J. Cha, and J. Park, ``{TOLEBI}: Learning fault-tolerant bipedal locomotion via online status estimation and fallibility rewards,'' arXiv preprint arXiv:2602.05596, 2026.

\bibitem{wu2026parkour}

Z. Wu \textit{et al.}, ``Perceptive humanoid parkour: Chaining dynamic human skills via motion matching,'' arXiv preprint arXiv:2602.15827, 2026.

\bibitem{domberg2026selfadapting}

F. Domberg and G. Schildbach, ``Self-adapting robotic agents through online continual reinforcement learning with world model feedback,'' arXiv preprint arXiv:2603.04029, 2026.

\bibitem{jayasinghe2026residual}

N. Jayasinghe \textit{et al.}, ``Residual control for fast recovery from dynamics shifts,'' arXiv preprint arXiv:2603.07775, 2026.

\bibitem{hu2024egocentric}

Y. Hu, B. Chen, and H. Lipson, ``Egocentric visual self-modeling for autonomous robot dynamics prediction and adaptation,'' arXiv preprint arXiv:2207.03386, 2022.

\bibitem{schulze2024high}

L. Schulze and H. Lipson, ``High-degrees-of-freedom dynamic neural fields for robot self-modeling and motion planning,'' in \textit{Proc. IEEE Int. Conf. Robot. Autom. (ICRA)}, 2024.

\end{thebibliography}

\end{document}
